% Sample LaTeX file for creating a paper in the Morgan Kaufmannn two
% column, 8 1/2 by 11 inch proceedings format.

\documentclass[letterpaper]{article}
\usepackage{uai2019}
\usepackage[margin=1in]{geometry}

\usepackage{microtype}
\usepackage{graphicx}
\usepackage{subfigure}
\usepackage{booktabs} % for professional tables

% hyperref makes hyperlinks in the resulting PDF.
% If your build breaks (sometimes temporarily if a hyperlink spans a page)
% please comment out the following usepackage line and replace
% \usepackage{icml2019} with \usepackage[nohyperref]{icml2019} above.
\usepackage{hyperref}

% Attempt to make hyperref and algorithmic work together better:
\newcommand{\theHalgorithm}{\arabic{algorithm}}

\usepackage{amsfonts}

\usepackage{amsmath}
\usepackage{amsthm}
\usepackage{bm}
\usepackage{color}
\usepackage{bbm}
\usepackage{color,graphicx}
%\usepackage[colorlinks]{hyperref}
\usepackage{natbib}
%\usepackage[numbers,sort]{natbib}
%\graphicspath{{../figures/}}
\usepackage{enumerate}
\usepackage{multirow}
\usepackage{cases}

\usepackage{comment}

\newcommand{\R}{\mathbb{R}}
\newcommand{\D}{\mathbb{D}}
\newcommand{\Z}{\mathbb{Z}}
\newcommand{\E}{\mathbb{E}}
\newcommand{\Prob}{\mathbb{P}}
\providecommand{\abs}[1]{\lvert#1\rvert}
\providecommand{\norm}[1]{\lVert#1\rVert}
\providecommand{\diag}{{\rm diag}}
\providecommand{\argmax}{{\rm argmax}}
\providecommand{\argmin}{{\rm argmin}}
\newcommand{\ds}{\displaystyle }
\usepackage{multicol}
\usepackage{subfigure}

%\bibliographystyle{abbrvnat}

\graphicspath{{../figures/}}

\newcommand{\cost}{\text{cost}}
\newcommand{\one}{1}
%\newcommand{\one}{1_m}
%\newcommand{\one}{\vec{1}}
%\newcommand{\one}{\bf{1}}
\newcommand{\rs}[1]{{\mbox{\scriptsize \sc #1}}}
\newcommand{\vc}[1]{{\boldsymbol #1}}
\newcommand{\vcn}[1]{{\bf #1}}
\newcommand{\vcs}[1]{\mbox{\footnotesize{\boldmath $#1$}}}
\newcommand{\sr}[1]{{\cal #1}}
\newcommand{\dd}[1]{\mathbb{#1}}
\newcommand{\ray}{{\rm r}}
\newcommand{\cp}{{\rm c}}
\newcommand{\br}[1]{\left\langle #1 \right\rangle}
\newcommand{\brb}[1]{\big\langle #1 \big\rangle}
\newcommand{\ol}{\overline}
\newcommand{\ul}{\underline}
\newcommand{\bx}{\bm{x}}
\newcommand{\by}{\bm{y}}
\newcommand{\vv}[1]{\bm{x_1}, \ldots, \bm{x_{#1}}}
\newcommand{\x}[1]{x^{(#1)}}
\newcommand{\z}[1]{z^{(#1)}}
\newcommand{\bmu}[1]{\bm{\mu}^{(#1)}}
\newcommand{\y}[1]{y^{(#1)}}
\newcommand{\m}[1]{\bm{m}^{(#1)}} % mean function
\newcommand{\Sig}[1]{\bm{\Sigma}^{(#1)}} % covariance function
\newcommand{\EI}{\mathrm{EI}}
\newcommand{\qEI}{\text{\it EI}}
\newcommand{\qEIg}{\text{\it d-EI}}
\newcommand{\VOIE}{\text{VOI}^\emptyset}
\newcommand{\qKGg}{\text{\it d-KG}}
\newcommand{\taKG}{\text{taKG}}
\newcommand{\taKGE}{\text{taKG}^\emptyset}
\newcommand{\qtaKG}{\text{q-taKG}}
\newcommand{\qtaKGE}{\text{q-taKG}^\emptyset}
\newcommand{\cfKG}{\text{cfKG}}
\newcommand{\qcfKG}{\text{q-cfKG}}

\newcommand{\A}{\mathbb{A}} % domain

\newtheorem{lemma}{Lemma}
\newtheorem{example}{Example}
\newtheorem{proposition}{Proposition}
\newtheorem{corollary}{Corollary}
\newtheorem{theorem}{Theorem}
%\theoremstyle{definition}
\numberwithin{equation}{section}
\newtheorem{definition}{Definition}
\newtheorem{remark}{Remark}

\newcommand{\red}[1]{{\color{red} #1}}
\newcommand{\blue}[1]{{\color{blue} #1}}
\newcommand{\green}[1]{{\color{green} #1}}
\newcommand{\pfcomment}[1]{{\color{blue} PF: #1}}
\newcommand{\pfedit}[1]{{\color{blue} PF: #1}}

% I want to use \S to refer to \mathcal{S} because I use it a lot,
% but I also want to be able to use the section symbol.
% This changes the command for the section symbol to \Section
\let\Section\S
\renewcommand{\S}{\mathcal{S}}




\newcommand{\eq}[1]{(\ref{eq:#1})}
\newcommand{\eqn}[1]{(\ref{eqn:#1})}
\newcommand{\lem}[1]{Lemma~\ref{lem:#1}}
\newcommand{\cor}[1]{Corollary~\ref{cor:#1}}
\newcommand{\thr}[1]{Theorem~\ref{thr:#1}}
\newcommand{\con}[1]{Conjecture~\ref{con:#1}}
\newcommand{\pro}[1]{Proposition~\ref{pro:#1}}
\newcommand{\que}[1]{Question~\ref{que:#1}}
\newcommand{\exa}[1]{Example~\ref{exa:#1}}
\newcommand{\ass}[1]{Assumption~\ref{ass:#1}}
\newcommand{\dfn}[1]{Definition~\ref{dfn:#1}}
\newcommand{\rem}[1]{Remark~\ref{rem:#1}}
\newcommand{\fig}[1]{Figure~\ref{fig:#1}}
\newcommand{\app}[1]{Appendix~\ref{app:#1}}
\newcommand{\sectn}[1]{\Section\ref{#1}}


\newcommand{\lemt}[1]{\ref{lem:#1}}
\newcommand{\cort}[1]{\ref{cor:#1}}
\newcommand{\thrt}[1]{\ref{thr:#1}}
\newcommand{\cont}[1]{\ref{con:#1}}
\newcommand{\prot}[1]{\ref{pro:#1}}
\newcommand{\exat}[1]{\ref{exa:#1}}
\newcommand{\dent}[1]{\ref{den:#1}}
\newcommand{\asst}[1]{\ref{ass:#1}}
\newcommand{\remt}[1]{\ref{rem:#1}}
\newcommand{\figt}[1]{\ref{fig:#1}}
\newcommand{\appt}[1]{\ref{app:#1}}
\newcommand{\sect}[1]{\ref{sect:#1}}

\newenvironment{mylist}[1]{\begin{list}{}
{\setlength{\itemindent}{#1mm}}
%{\setlength{\itemsep}{0ex plus 0.2ex}}
{\setlength{\itemsep}{0pt}}
{\setlength{\parsep}{0.5ex plus 0.2ex}}
\setlength\bibitemsep{0pt}
{\setlength{\labelwidth}{10mm}}
}{\end{list}}

\newcommand{\stedit}[1]{{\color{blue} #1}}
\newcommand{\stcomment}[1]{{\color{blue} ST: #1}}

\newcommand{\Y}{\mathcal{Y}}
\newcommand{\loss}{L}

\usepackage{color}
\usepackage[dvipsnames]{xcolor}
\newcommand{\xxcomment}[4]{\textcolor{#1}{[$^{\textsc{#2}}_{\textsc{#3}}$ #4]}}

\newcommand{\agw}[1]{\xxcomment{red}{A}{W}{#1}}
\newcommand{\pf}[1]{\xxcomment{blue}{P}{F}{#1}}


% Set the typeface to Times Roman
\usepackage{times}



\title{Practical Multi-fidelity Bayesian Optimization for Hyperparameter Tuning}
% Practical multi-fidelity Bayesian optimization of iterative machine learning algorithms
% Practical multi-fidelity Bayesian optimization for hyperparameter tuning

\author{
Jian Wu \quad
Saul Toscano-Palmerin \quad
Peter I. Frazier \quad
Andrew Gordon Wilson \\
Cornell University
}

% The author names and affiliations should appear only in the accepted paper.
%
%\author{ {\bf Harry Q.~Bovik\thanks{Footnote for author to give an
%alternate address.}} \\
%Computer Science Dept. \\
%Cranberry University\\
%Pittsburgh, PA 15213 \\
%\And
%{\bf Coauthor}  \\
%Affiliation          \\
%Address \\
%\And
%{\bf Coauthor}   \\
%Affiliation \\
%Address    \\
%(if needed)\\
%}

\begin{document}

\maketitle

\begin{abstract}

Bayesian optimization is popular for optimizing time-consuming black-box objectives.
%\pf{including hyperparameter tuning of large machine models}. 
Nonetheless, for hyperparameter tuning in deep neural networks,
the time required to evaluate the validation error for even a few hyperparameter settings remains a bottleneck.
Multi-fidelity optimization promises relief using cheaper proxies to such objectives ---  for example, validation error for a network trained using a subset of the training points or fewer iterations than required for convergence.
We propose a highly flexible and practical approach to multi-fidelity Bayesian optimization, 
focused on efficiently optimizing hyperparameters for iteratively trained supervised learning models.
We introduce a new acquisition function, the trace-aware knowledge-gradient, which efficiently leverages both multiple continuous fidelity controls and trace observations --- values of the objective at a sequence of fidelities, available when varying fidelity using training iterations. 
We provide a provably convergent method for optimizing our acquisition function
and show it outperforms state-of-the-art alternatives for  hyperparameter tuning of deep neural networks and large-scale kernel learning. 

%\agw{This abstract is way too long}
%We consider multi-fidelity Bayesian optimization with ``trace" observations \agw{it's not clear why we're calling these ``trace'' observations}: observing at one fidelity also provides observations at lower fidelities. This occurs when tuning hyperparameters to optimize validation error of a deep neural network using the number of training iterations as a fidelity: after computing validation error at $s$ training iterations, we can compute validation error at $s'<s$ training iterations essentially for free.  
%We propose a new multi-fidelity Bayesian optimization method that effectively utilizes traces, making four key innovations.
%(1) We support multiple simultaneous fidelity controls, including one or more that provide a trace.
%When tuning deep neural networks, this lets our method control both training iterations and training data to find good solutions more quickly, while previous methods allow only one fidelity control.
%(2) We develop a novel parameter-free method for overcoming the tendency common among multi-fidelity methods to sample too frequently at fidelities that provide almost no information at almost no cost. The method works both with and without traces.
%(3) We save computation by including in inference only the most valuable parts of the trace.
%(4) We leverage the ability to continue a previously initiated trace at reduced cost.
%Numerical experiments show that our method outperforms state-of-the-art algorithms when optimizing synthetic functions, tuning feedforward neural networks on MNIST, tuning convolutional neural networks on CIFAR-10 and SVHN, and in large-scale kernel learning.
\end{abstract}

\section{INTRODUCTION}
In hyperparameter tuning of machine learning models, we seek to find hyperparameters $x$ in some set ${A} \subseteq \mathbb{R}^d$ to minimize the validation error $f(x)$, i.e., to solve
\begin{equation}
\min_{x\in\mathbb{A}} f(x)
\label{eqn:min_f}
\end{equation}
Evaluating $f(x)$ can take substantial time and computational power \citep{bergstra2012random}
and may not provide gradient evaluations. Bayesian optimization, which requires relatively few 
function evaluations, provides a compelling approach to such optimization problems \citep{jones1998efficient, snoek2012practical}.

As the computational expense of training and testing a modern deep neural network for a single set of hyperparameters has grown, researchers have sought to supplant some evaluations of $f(x)$ with computationally inexpensive low-fidelity approximations.  
Conceputally, an algorithm can use low-fidelity evaluations to quickly identify a smaller set of promising hyperparameter settings, and then later focus more expensive high-fidelity evaluations within this set to refine its estimates.
% Leveraging evaluations of multiple fidelities promises to make hyperparameter tuning substantially more efficient.

Pioneering multi-fidelity approaches focused on hyperparameter tuning for deep neural networks
include the Bayesian optimization methods FaBOLAS \citep{klein2016fast,klein2015towards}, Freeze-Thaw Bayesian Optimization \citep{swersky2014freeze}, BOCA \citep{kandasamy2017multi}, predictive entropy search for a single continuous fidelity \citep{mcleod2017practical}, early-stopping SMAC \citep{domhan2015speeding}, and
Hyperband \citep{li2016hyperband}.
This work builds on
%Interest in such approaches for hyperparameter tuning draws from 
earlier multi-fidelity optimization approaches \citep{huang2006sequential,lam2015multifidelity,poloczek2016multi} focused on low-fidelity approximation of physics-based computational models.


These validation error approximations 
% appropriate for multifidelity hyperparameter tuning
perform the same training and testing steps as in standard Bayesian optimization, but control fidelity with fewer training iterations than required for convergence, fewer training data points, or fewer validation data points. 
These approximations present unique opportunities not typically considered in the multifidelity literature, even within the portion focused on hyperparameter tuning. 
First, we observe a full \emph{trace} of performance with respect to training iterations, rather than just a single performance value at the chosen fidelity.
Indeed, training with $s$ iterations produces evaluations of the low-fidelity approximation for {\it all} training iterations less than or equal to $s$.  
% Using more of this information is an opportunity to improve performance.
Second, by caching state after completing $s$ iterations, we can significantly reduce computation time when later evaluating for $s'>s$ evaluations.  This allows quickly evaluating low-fidelity approximations to the validation error for many hyperparameter settings, then later returning to those most promising hyperparameter settings to cheaply obtain more accurate observations.
Third, we may simultaneously alter fidelity along several continuous dimensions (iterations, training data, validation data), rather than modifying one continuous fidelity control or choosing from among a discrete set of ambiguously related fidelities.

%Despite its promise, observing traces over training iterations raises key challenges:
%(1) Balancing between increasing training iterations at previously evaluated hyperparameters vs. starting from scratch at new hyperparameters;
%(2) Extracting the most useful information from a set of training iterations without burdening a Gaussian process inference procedure (which scales as the cube of the number of observations) with many observations. 

In this paper, we propose the trace-aware knowledge gradient (taKG) for Bayesian optimization with multiple fidelities. taKG is distinctive in that it leverages both trace information and multiple fidelity controls at once, efficiently selecting training size, validation size, number of training iterations, and hyperparameters to optimize. Moreover, we provide a provably-convergent method for maximizing this acquisition function.
taKG addresses the challenges presented by trace observations by considering the reduced cost of adding iterations at a previously evaluated point, and using an intelligent selection scheme to choose a subset of the observed training iterations to include in inference.
Additionally, taKG can be used in either a batch or sequential setting, and can also efficiently leverage gradient information if it is available.

We present two variants of our trace-aware knowledge-gradient acquisition function, one for when the cost of sampling is substantial over the whole fidelity space (even when using few training points or iterations), and the other for when the cost and value of information vanish as fidelity decreases to 0.  The first form we refer to simply as taKG, and the second as 0-avoiding taKG ($\taKGE$) because it avoids the tendency of multi-fidelity other methods to measure repeatedly at near-0 fidelities even when these low fidelities provide almost no useful information.  Alternative approaches \citep{mcleod2017practical,klein2016fast} add and tune a fixed cost per sample to avoid this issue, while $\taKGE$ does not require tuning. 

Furthermore, we present a novel efficient method to optimize these acquisition functions, even though they cannot be evaluated in closed form.
 This method first constructs a stochastic gradient estimator which it then uses within multistart stochastic gradient ascent. We show that our stochastic gradient estimator is unbiased and thus asymptotically consistent, and the resulting stochastic gradient ascent procedure converges to a local stationary point of the acquisition function.
 
Our numerical experiments demonstrate a significant improvement over state-of-the-art alternatives, including FaBOLAS \citep{klein2016fast,klein2015towards}, Hyperband \citep{li2016hyperband}, and 
BOCA \citep{kandasamy2017multi}. Our approach is also applicable to problems that do not have trace observations, but use continuous fidelity controls, and we additionally show strong performance in this setting.

In general, efficient and flexible multifidelity optimization is of crucial practical importance, as evidenced by growing momentum in this research area. Although Bayesian optimization has shown great promise for tuning hyperparameters of machine learning algorithms, computational bottlenecks have remained a major deterrent to mainstream adoption. With taKG, we leverage crucial trace information, while simultaneously providing support for several fidelity controls, providing remarkably efficient optimization of expensive objectives. This work is intended as a step towards the renewed practical adoption of Bayesian optimization for machine learning.

% (3) Modifying our decisions about what hyperparameters and fidelities to evaluate in light of the fact that we will observe not just a single low-fidelity approximation but a range indexed by the training iterations performed.

\begin{comment}
\agw{should there be a related work section? the following paragraphs, up to contributions, (1) seem too long-winded and out of place for an introduction; (2) way too specialized; (3) defensive and may alienate many of our potential reviewers}
\citet{li2016hyperband} and \citet{domhan2015speeding} avoid these challenges by placing artificial limits on starting and stopping evaluation of hyperparameters: \citet{li2016hyperband} simply never adds new promising hyperparameters to its set of candidates;
while \citet{domhan2015speeding} fully evaluates all but extremely poorly performing candidates, and never restarts hyperparameter evaluation after stopping it.
In addition, \citet{li2016hyperband} can consider only one fidelity at a time (either training iterations or training data, but not both) and \citet{domhan2015speeding} cannot consider training data.  These artificial limits simplify analyses, but reduce potential benefits. \agw{I think the way this paragraph is phrased would upset these authors}

\citet{kandasamy2017multi} avoids these challenges by ignoring the opportunity offered by the unique characteristics of training iterations as a fidelity, instead adopting a model closer to traditional multi-fidelity Bayesian optimization: it supposes we do not observe data from training iterations smaller than the one requested, which discards information about the training loss' rate of decrease; and it imagines evaluating $s$ iterations does not reduce the cost of evaluating $s'>s$ iterations in the future.

\citet{klein2016fast} and \citet{mcleod2017practical} avoid these challenges by not considering training iterations to be a fidelity. Instead, they focus on training data size alone, where these challenges do not arise.
While this simplifies analysis, 
we argue that not considering training data size and training iterations together hampers the speed at which high-quality solutions may be found.
% despite the fact that both are important controls on training cost and controlling them together would seem to offer a significant opportunity to improve performance.

\citet{swersky2014freeze}, an unpublished working paper \agw{shouldn't say this}, considers using training iterations as a fidelity most directly, but does not allow inclusion of other fidelities (e.g., training data size) and its approach does not extend easily in that direction.  It also makes approximations that may harm performance:
for example, 
the set of new hyperparameters considered for evaluation with predictive entropy search is limited to a basket of size three.
% pre-convergence training loss across different sets of hyperparameters is assumed conditionally independent given the limiting value;
\citet{domhan2015speeding} cites a personal communication with the authors reporting the method does not perform well on tuning deep networks. \agw{this paragraph seems quite awkward}

% \cite{li2016hyperband} avoids challenge (1) by using a tournament-style approach that does not allow starting at new sets of hyperparameters, but this prevents it from leveraging revealed from its initial evaluations to choose a promising new set of hyperparameters at which to evaluate.  It avoids challenge (2) by ignoring all but the most recent iteration, but this hurts its ability to incorporate trends into its projection for a set of hyperparameters' asymptotic performance.  \cite{domhan2015speeding} avoids challenge (1) by choosing not to restart evaluation at a point once it has been terminated, which hampers its ability to get a rough approximation of the response surface through evaluations at many points 

% A third challenge arises because an evaluation with less than the full training iterations or training data does not result in evaluating the objective.  This makes it harder to define a measure of expected improvement as one's acquisition function within a principalled framework.  While \cite{domhan2015speeding} continues to use expected improvement as a heuristic, the approaches \cite{swersky2014freeze} and \cite{klein2016fast} instead propose methods based on entropy search.  This in turn reveals its own challenges, because the ideal entropy search strategy is hard to realize computationally, and requires approximations whose impact on solution quality is as yet unclear.

% Perhaps because of these challenges, we are unaware of a previously proposed Bayesian optimization method capable of simultaneously choosing both training data size and training iterations, despite the fact that both are important controls on training cost and controlling them together would seem to offer a significant opportunity to improve performance.
%\vskip -1in


\paragraph{Contributions:}
\agw{theoretical contributions?}
We propose the trace-aware knowledge gradient, which can select the training data size, training iterations, and hyperparameters at which to evaluate.  Moreover, we provide a provably-convergent method for maximizing this acquisition function.
Our approach is applicable to problems with trace observations along one or more fidelity controls, including hyperparameter optimization, while varying both training iterations and training data \agw{isn't that what you just said, two sentences ago?}.
It addresses the challenges presented by trace observations by considering the reduced cost of adding iterations at a previously evaluated point, and using an intelligent selection scheme to choose a subset of the observed training iterations to include in inference.
It can be used in either a batch or sequential setting, and can leverage gradient information if it is available.
\agw{mention theoretical contributions?}


% Maximizing taKG selects the point or set of points to evaluate next that maximizes the ratio of the value of information from evaluation against its cost.  We provide a novel stochastic gradient descent approach that can be used to maximize the taKG, provide a convergence result showing that this SGD approach indeed converges to the maximizer of the taKG.  
%We also provide a proof of consistenty showing that in the limit as the number of evaluations we do grows large, our estimator of the optimal set of hyperparameters converges to the correct global optimum.

We present two variants of our trace-aware knowledge-gradient acquisition function, one for use when the cost of sampling is substantial over the whole fidelity space \agw{e.g., for any number of training points}, and the other for when the cost and value of information vanish as fidelity decreases to 0.  The first form we refer to simply as taKG, and the second as 0-avoiding taKG ($\taKGE$) because it avoids the tendency of other methods to measure repeatedly at near-0 fidelities even when these low fidelities provide almost no useful information.  Other methods \citep{mcleod2017practical,klein2016fast} add and tune a fixed cost per sample to avoid this issue, while $\taKGE$ does not require tuning. 

While these acquisition functions cannot be evaluated in closed form, presenting a challenge to their optimization and thus use within a Bayesian optimization framework, we present a novel efficient computational method that optimizes them despite this challenge.  This method first constructs a stochastic gradient estimator which it then uses within multistart stochastic gradient ascent.  Justifying use, we show that our stochastic gradient estimator is unbiased and thus asymptotically consistent, and the resulting stochastic gradient ascent procedure converges to a local stationary point of the acquisition function.

Our numerical experiments demonstrate a significant improvement over FaBOLAS, Hyperband, and BOCA.  We do not compare against \citet{swersky2014freeze} as the method is incompletely described (e.g., how the 2$^{\mathrm{nd}}$ and 3$^{\mathrm{rd}}$ candidate solutions are added to the basket) and code is not publicly available.\agw{if we include comments like this, they should be as a footnote in the experiments section}.

%In contrast with \cite{domhan2015speeding} and \cite{li2016hyperband} our approach allows itself the full range of actions: starting, stopping, and continuing evaluations at sets of hyperparameters as it becomes productive to do so without limit.  

%In contrast with \citet{kandasamy2017multi} and \citet{klein2016fast} our approach considers the unique characteristics of training iterations as a fidelity, leveraging the ability to increase training iterations at a previously evaluated point at a lower cost, and selecting a .  

%In contrast with \cite{swersky2014freeze} our approach includes an 
% Our conceptual approach provides a Bayes-optimal batch of measurements to perform when only one batch remains.  Our theoretical analysis shows that the error in implementing this conceptual approach vanishes as computing power increases, in contrast with existing work on entropy search methods \cite{swersky2014freeze,klein2016fast} that make approximations that do not clearly vanish as computation grows.

% Our approach leverages existing work on knowledge-gradient acqusition functions \cite{wu2017bayesian}, and builds novel approaches for selecting which portions of the trace of training iterations to include to both extract the most meaningful data and control inference overhead. 

Our approach is also applicable to problems that do not have trace observations, but use continuous fidelity controls.
Indeed, other than BOCA and the proposed taKG, we are unaware of other methods that explicitly treat more than one continuous fidelity \agw{meaning, e.g., training points and number of iterations?}, and taKG outperforms BOCA on problems without trace observations in our experiments.
This may be because BOCA uses a two-stage process that selects the point to evaluate without considering the fidelity, and only selects the fidelity afterward.  In contrast, taKG selects the point and fidelity jointly, and considers the impact of fidelity choice on the best point to evaluate.
taKG is also the first algorithm to allow both continuous fidelities and batch or derivative observations within one framework.

%% saving space
%While these joint capabilities appear to be new, it may be possible to adapt other batch or derivative-enabled methods to the continuous fidelity setting, and to adapt other single-control methods to include multiple controls.  Thus, we emphasize that our primary metholodical contributions are taKG's ability to leverage the specific nature of trace observations and its robustness to a vanishing cost and value of information.


\paragraph{Other related work:}
\agw{We can probably incorporate some of these citations into our descriptions of related work, but don't require this paragraph}.
Beyond related work discussed above proposing other approaches to multi-fidelity Bayesian optimization in settings related to the one with trace observations we consider here, our work is also related methodologically to work on acquisition functions in non-trace settings.   First, our acquisition function is a knowledge-gradient acquisition function, and is thus related to work on knowledge-gradient acquisition functions for other problem settings \citep{frazier2009knowledge,wu2016parallel}.  Second, the technique we propose for optimizing our acquisition function generalizes a recently proposed computational method for single-fidelity optimization in \citet{wu2017bayesian} to trace observations.  Third, our acquisition functions is given by a ratio of a value of information to the cost of acquiring that information. Acquisition functions based on ratios of value to cost have been studied in other contexts \citep{klein2016fast,mcleod2017practical,poloczek2016multi}.


\end{comment}

\section{THE taKG AND $\taKGE$ ACQUISTION FUNCTIONS}
\label{sect:taKG}
In this section we define the trace-aware knowledge-gradient acquisition function.
\sectn{sec:obj} defines our formulation of multi-fidelity optimization with traces and continuous fidelities, along with our inference procedure.
\sectn{sect:valuing} describes a measure of expected solution quality possible after observing a collection of fidelities within a trace.
\sectn{sect:voi} uses this measure to define the $\taKG$ acquisition function,
and \sectn{sect:mvoi} defines an improved version, $\taKGE$, appropriate for settings in which the 
the cost and value of information vanish together (for example, as the number of training iterations declines to 0).
\sectn{sect:optimize} then presents a computational approach for maximizing the $\taKG$ and $\taKGE$ acquisition functions and theoretical results justifying its use.  \sectn{sect:warm-start} discusses warm-starting previously stopped traces, and \sectn{sect:extensions} briefly discusses generalizations to batch and derivative observations.

% , which values sampling a point at a fidelity according to the ratio of the value of the information gained to the cost of doing so.  This method works well when the sampling cost stays substantially away from 0 for all fidelities, but is sensitive to model misspecification when the cost and value of information can vanish together (for example, as the number of training iterations declines to 0).  We discuss this issue in \sectn{sect:mvoi} and describe a variation of our acquisition function, the 0-avoiding trace-aware knowledge gradient ($\taKGE$), which overcomes this challenge.

% This approach is based on an envelope-theorem based method for obtaining unbiased samples of the gradient of the acquisition function. 

% propose the trace-aware knowledge gradient ($\taKG$) and the 0-avoiding trace-aware knowledge gradient ($\taKGE$), a pair of novel Bayesian optimization algorithm that exploit inexpensive low-fidelity approximations, of which one fidelity control may optionally provide trace observations, such as the number of training iterations.

\subsection{Problem Setting}
\label{sec:obj}
We model our objective function and its inexpensive approximations by a real-valued function $g(x, s)$ where our objective is $f(x) := g(x, 1)$ and $s \in [0,1]^m$ denotes the $m$ fidelity-control parameters.   (Here, $1$ in $g(x,1)$ is a vector of $1$s.)  
We assume that our fidelity controls have been re-scaled so that $1$ is the highest fidelity and $0$ the lowest.  
$g(x,s)$ can be evaluated, optionally with noise, at a cost depending on $x$ and $s$. 

We let $B(s)$ be the additional fidelities observed for free when observing fidelity $s$.  Although our framework can be easily generalized, we assume that $B(s)$ is a cross product of sets of the form either $[0,s_i]$ (\emph{trace fidelities}) or $\{s_i\}$ (\emph{non-trace fidelities}). 
We let $m_1$ denote the number of trace fidelities and $m_2$ the number of non-trace fidelities.
We also assume that the cost of evaluation is non-decreasing in each component of the fidelity.

For example, consider hyperparameter tuning with $m=2$ fidelities: first is the number of training iterations; % expressed as the fraction of some maximum number of training iterations;
second is the amount of training data.  %, expressed as a fraction of some maximum amount.  
Each is bounded between $0$ and some maximum value, and $s_i \in [0,1]$ specifies training iterations or number of training data points as a fraction of this maximum value.
Then, $B(s) = [0,s_1] \times \{s_2\}$, because we observe results for the number of training iterations ranging from $0$ up to the number evaluated.  
% If we had a third fidelity that did not provide trace observations, $B(s) = [0,s_1] \times \{s_2\} \times \{s_3\}$.  
If the amount of validation data is another trace fidelity, we would have: $B(s) = [0,s_1] \times \{s_2\} \times [0,s_3]$. 

We model $g$ using Gaussian process regression jointly over $x$ and $s$, assuming that observations are perturbed by independent normally distributed noise with mean $0$ and variance $\sigma^2$.  Each evaluation consists of $x$, a vector of fidelities $s$, and a noisy observation of $g(x,s')$ for each fidelity $s'$ in $B(s)$.  For computational tractability, in our inference, we will choose to retain and incorporate observations only from a subset $\mathcal{S} \subseteq B(s)$ of these fidelities with each observation.  After $n$ such evaluations, we will have a posterior distribution on $g$ that will also be a Gaussian process, and whose mean and kernel we refer to by $\mu_n$ and $K_n$.  We describe this inference framework in more detail in the supplement.

We model the logarithm of the cost of evaluating $g$ using a separate Gaussian process, updated after each evaluation,
and let $\cost_n(x,s)$ be the predicted cost after $n$ evaluations.
We assume for now that the cost of evaluation does not depend on previous evaluations, and then discuss later in \sectn{sect:warm-start} an extension to warm-starting evaluation at higher fidelities using past lower-fidelity evaluations.



\subsection{Valuing Trace Observations}
\label{sect:valuing}

Before defining the $\taKG$ and $\taKGE$ acquisition functions, we define a function $\loss_n$ that quantifies the extent to which observing trace information improves the quality of our solution to \eqref{eqn:min_f}.

Let $\E_n$ indicate the expectation with respect to the posterior $\mathbb{P}_n$ after $n$ evaluations.  Given any $x$ and set of fidelities $\S$, 
we will define a function $\loss_n(x,\S)$ to be the expected loss (with respect to the time-$n$ posterior) of our final solution to \eqref{eqn:min_f} if we are allowed to first observe $x$ at all fidelities in $\S$.

To define this more formally, let $\Y(x,\S)$ be a random vector comprised of observations of $g(x,s)$ for all $s \in \S$.  Then, the conditional expected loss from choosing a solution $x'$ to \eqref{eqn:min_f} after this observation is $\E_n\left[g(x',1) \mid \Y(x,\S)\right]$.  
This quantity is a function of $x$, $\S$, $\Y(x,\S)$, and the first $n$ evaluations, and 
can be computed explicitly using formulas from GP regression given in the supplement.

We would choose the solution for which this is minimized, giving a conditional expected loss of $\min_{x'} \E_n\left[g(x',1) \mid \Y(x,\S) \right]$. This is a random variable under the time-$n$ posterior whose value depends on $\Y(x,\S)$.
We finally take the expected value under the time-$n$ posterior to obtain $\loss_n(x,\S)$:
\begin{equation} 
\begin{split}
&\loss_n(x,\S) := \E_n\left[ \min_{x'} \E_n\left[g(x',1) \mid \Y(x,\S) \right]\right] \\
&\!=\!\int\!\mathbb{P}_n\!\left(\Y(x,\S)\!=\!y\right) \min_{x'} \E_n\left[g(\!x'\!,\!1\!)\!\mid\!\Y(\!x\!,\!\S)\!=\!y \right]\, dy,
\end{split}
\end{equation}
where the integral is over all $y \in \mathbb{R}^{|\S|}$.

We compute this quantity using simulation.  To create one replication of this simulation we first simulate $(g(x,s) : s \in S)$ from the time-$n$ posterior.  We then add simulated noise to this quantity to obtain a simulated value of $\Y(x,\S)$.  We then update our posterior distribution on $g$ using this simulated data, allowing us to compute $\E_n\left[g(x',1) \mid \Y(x,\S)\right]$ for any given $x'$ as a predicted value from GP regression.
We then use continuous optimization method designed for inexpensive evaluations with gradients (e.g., multi-start L-BFGS) to optimize this value, giving one replication of $\min_{x'} \E_n\left[g(x',1) \mid \Y(x,\S)\right]$.  We then average many replications to give an unbiased and asymptotically consistent estimate of $\loss_n(x,\S)$.

In a slight abuse of notation, we also define $\loss_n(\emptyset) = \min_{x'} \E_n\left[g(x',1) \right]$.  This is the minimum expected loss we could achieve by selecting a solution without observing any additional information.
This is equal to $\loss_n(x,\emptyset)$ for any $x$.

The need to compute $\loss_n(x,\S)$ via simulation will present a challenge when optimizing acquisition functions defined in terms of it.  Below, in \sectn{sect:optimize} we will overcome this challenge via a novel method for simulating unbiased estimators of the gradient of $\loss_n(x,\S)$ with respect to $x$ and the components of $\S$.  First, however, we define the $\taKG$ and $\taKGE$ acqisition functions.

\subsection{Trace-aware Knowledge Gradient (taKG)}
\label{sect:voi}

%Here, we consider the general continuous-fidelity setting with trace observations.
%Our acquisition function is built with the assumption that we will build information from the trace into our statistical model.  However, naively including all observations from the trace into the statistical model is computationally prohibitive with a cost of the order $\mathcal{O}(N^3T^3)$ where $N$ is the number of configurations and $T$ is the number of points available on the trace. We instead retain two observations from each trace, with the intuition that adding a second point provides additional information beyond the first about the rate at which the test loss is decreasing.  We choose not to include more than two to minimize time spent in inference, and because adding additional points provides decreasing marginal returns.  Our acquisition function will guide which additional point to include beyond the final point, and will modify its decisions beyond a trace-free knowledge-gradient approach in light of the fact that we obtain this extra information.


The $\taKG$ acquisition function will value observations of a point and a collection of fidelities according to the ratio of the reduction in expected loss (as measured using $\loss_n$) that it induces, to its computational cost.

While evaluating $x$ at a fidelity $s$ in principle provides observations of $g(x,s')$ at all $s' \in  B(s)$, we choose to retain and include in our inference only a subset of the observed fidelities $\S \subseteq B(s)$.  This reduces computational overhead in GP regression.  In our numerical experiments, we take the cardinality of $\S$ to be either 2 or 3, though the approach also allows increased cardinality. 

With this in mind, the $\taKG$ acquisition function at a point $x$ and set of fidelities $\S$ at time $n$ is
\begin{equation*}
\taKG_n(x, \mathcal{S}) := \frac{\loss_n(\emptyset) - \loss_n(x,\S)}{\cost_n(x, \max \S)},
\end{equation*}
where we also refer to the numerator as the {\it value of information} \citep{Ho66},
$\text{VOI}_n(x,\S) := \loss_n(\emptyset) - \loss_n(x,\S)$.
Thus, $\taKG$ quantifies the value of information per unit cost of sampling.

The cost of observing at all fidelities in $\S$ is taken here to be the cost of evaluating $g$ at a vector of fidelities equal to the elementwise maximum, $\max \mathcal{S} := (\max_{s \in \mathcal{S}} s_i: 1 \le i \le m)$.  
%Under the assumption that the cost is non-decreasing componentwise in the fidelity, 
This is the least expensive fidelity at which we could observe $\S$. 
%\pf{would be worth adding an example if we have time.}

%\begin{figure}[tb]
%\centering
  %\subfigure
  %\centering
  %\includegraphics[width=140pt, height = 100pt]{demo}
%\vspace{-8pt}
%\caption{\small Illustration of the evaluated fidelity (black dot on the top right) and the retained set of fidelities $\mathcal{S}^{(n+1)}$ (black dots) in a problem with two fidelities, one of which provides a trace (training iterations) and the other of which does not (training data size).}
%\vspace{-12pt}
%\label{Fig_enhancing_set}
%\end{figure}

The taKG algorithm chooses to sample at the point $x$, fidelity $s$, and additional lower-fidelity point(s) $\S \setminus \{s\}$ to retain that jointly maximize the $\taKG$ acquisition function, among all fidelity sets $\S$ with limited cardinality $\ell$.

\begin{equation}
\max_{x,s,\mathcal{S}:\mathcal{S}\subseteq B\left(s\right),\left|\mathcal{S}\right|=\ell,s\in\mathcal{S}} \taKG_n \left(x, \mathcal{S} \right).
\label{eqn:max-taKG}
\end{equation}
This is a continuous optimization problem whose decision variable is described by $d + \ell m_1 + m_2$ real numbers. $d$ describe $x$, $m=m_1+m_2$ describe $s$, and $(\ell-1)m_1$ describe $\S \setminus \{s\}$.

%\subsection{Reducing sensitivity to model misspecification tendency to sample at fidelity $0$}
%\subsection{Sensitivity to model-misspecification at low-fidelities}
\subsection{0-avoiding Trace-aware Knowledge Gradient ($\taKGE$)}
\label{sect:mvoi}
The $\taKG$ acquisition function uses the value of information per unit cost of sampling.  
When the value of information and cost become small simultaneously, as when we shrink training iterations or training data to 0 in hyperparameter tuning, this ratio becomes sensitive to misspecification of the GP model on $g$. We first discuss this issue, and then develop a version of $\taKG$ for these settings.

To understand this issue, we first observe the value of information for sampling $g(x,s)$, for any $s$, is strictly positive when the kernel has strictly positive entries.
%as it does for the M\`atern and other widely-used kernels.  
\begin{proposition}
If the kernel function $K_n((x, s), (x', 1)) > 0$ for any $x, x' \in \mathbb{A}$, then for any $x \in \mathbb{A}$ and any $s\in[0,1]^m$, $\text{VOI}_n(x, \left\{ s\right\}) > 0$.
\label{prop:positive}
\end{proposition}

Proposition~\ref{prop:positive} holds even if $s=0$, or has some components set to 0.
Thus, if the estimated cost at such extremely low fidelities is small relative to the (strictly positive) value of information there, $\taKG$ may be drawn to sample them, even though the value of information is small.  We may even spend a substantial portion of our budget evaluating $g(x,0)$ at different $x$.
This is usually undesirable.

For example, in hyperparameter tuning with training iterations as our fidelity, 
fidelity $0$ corresponds to training a machine learning model with no training iterations.
This would return the validation error on initial model parameter estimates.
While this likely provides {\it some} information about the validation error of a fully trained model, 
specifying a kernel over $g$ that productively uses this information from a large number of hyperparameter sets $x$ 
would be challenging.

This issue becomes even more substantial when considering training iterations
and training data together, as we do here, because cost nearly vanishes
as either fidelity vanishes.  Thus, there are many
fidelities at each $x$ that we may be drawn to oversample.

This issue is not specific to \taKG. 
It also occurs in previous literature 
\citep{klein2016fast,mcleod2017practical,klein-bayesopt17}
when using the ratio of information gain to cost in 
an entropy search or predictive entropy search method based on the same predictive model.

% This issue can be partially addressed by adding a fixed cost to all evaluations, which might be motivated by the overhead of optimizing the acquisition function as part of the BO algorithm.  However, this cost might be small or of the same order of magnitude as evaluating at fidelity $0$.  Thus, even when including this fixed cost, a misspecified GP model may lead us to sample excessively at small fidelities.  

To deal with this issue, \citet{klein2016fast,mcleod2017practical} artificially inflate the cost of evaluating at fidelity $0$ to penalize low fidelity evaluations.  Similarly, \citet{klein-bayesopt17} recommends adding a fixed cost to all evaluations motivated by the overhead of optimizing the acquisition function, but then recommends setting this to the same order of magnitude as a full-fidelity evaluation even though the overhead associated with optimizing a BO acquisition function using well-written code and efficient methodology will usually be substantially smaller. As a result, any fixed cost must be tuned to the application setting to avoid oversampling at excelssively small fidelities while still allowing sampling at moderate fidelities.
%\pfcomment{Jian: is this updated paragraph a good representation of the literature on this?}

Here, we propose an alternate solution that we find works well without tuning,
focusing on the setting where the cost of evaluation becomes small as the smallest component in $s$ approaches 0.
% Our solution is designed specifically for 

% $\text{cost}(x,s) \approx 0$ when $\min(s)=0$.
% , as it does in tuning hyperparameters of machine learning models. 
% Our alternate solution modifies our value of information calculation by first constructing a 

We first define $C(s) = \cup_{i=1}^m \{ s' : s'_i = 0, s'_j = s_i\ \forall j\ne i \}$
to be the set of fidelities obtained by replacing one component of $s$ by $0$.
We then let $C(\S) = \cup_{s \in \S} C(s)$ be the union of these fidelities over $s \in \S$.
For example, suppose $s_1$ is a trace fidelity (say, training iterations), $s_2$ is a non-trace fidelity (say, training data size), and $\S = \{ (1/2, 1), (1,1) \}$, corresponding to an evaluation of $g$ at $s=(1,1)$ and retention of the point $(1/2,1)$ from the trace $B((1,1))$.
Then $C(\mathcal{S}) = \{ (0,1), (1/2, 0), (1,0) \}$.  %Fig.~\ref{fig:C} illustrates this.

%\begin{figure}[htb]
%\centering
  %\includegraphics[width=155pt, height = 95pt]{demo}
%\vspace{-8pt}
%\caption{\small The black dots are the fidelity we choose and the middle point we retain on the trace, i.e. $\mathcal{S}$, the blue dots are the corresponding $C(\mathcal{S})$. \label{fig:C}}
%\label{demo}
%\end{figure}

%When costs becomes small as the smallest component in $s$ approaches 0,
%In the settings where the issue discussed above arises,
Fidelities in $C(\S)$ (for any $\S$) are extremely inexpensive to evaluate and provide extremely small but strictly positive value of information.
These, and ones close to them, are ones we wish to avoid sampling, even when $\taKG$ is large.

To accomplish this, we modify our value of information 
$\text{VOI}_n(x,\S) = \loss_n(\emptyset) - \loss_n(x,\S)$
to suppose free observations $\Y(x,s')$ will be provided of these problematic low-fidelity $s'$.
  Our modified value of information will suppose these free observations will be provided to both the benchmark, previously set to $\loss_n(\emptyset)$, and to the reduced expected loss, previously set to $\loss_n(x,\S)$, achieved through observing $x$ at fidelities $\S$.
The resulting modified value of information is
\begin{equation*}
\VOIE_n(x,\S) = \loss_n(x,C(\S)) - \loss_n(x,\S \cup C(\S)) \notag \\
\end{equation*}
We emphasize our algorithm will not evaluate $g$ at fidelities in $C(\S)$.  Instead, it will simulate these evaluations according to the algorithm in \sectn{sect:valuing}.

%If $\text{cost}^{(n)}(x,s) \le \epsilon$ when $\min(s) = 0$ and is strictly greater than $\epsilon$ otherwise, then
%$C(\S)$ to be the set consisting of those fidelities obtained by replacing one component of one fidelity $s'\in\mathcal{S}$ by $0$.  

% This will be the set of fidelities $s'$ with near 0 cost whose corresponding $g(x,s')$ we will condition on when computing $\mu_*^{(n)}$.
% Let $g(x,C(\mathcal{S})) = \{ g(x,s') : s' \in C(\mathcal{S}) \}$ be the value of $g(x,\cdot)$ at all of the fidelities in $C(\mathcal{S})$.

% In our value of information calculation, we suppose that the benchmark value received $\mu_*^{(n)}$ when computing the VOI is already conditioned on $g(x,s')$ for some fidelities with cost near 0.  Indeed, if sampling these fidelities has $0$ cost, we could obtain their values for free.  In the algorithm, we do not evaluate a model with fidelity 0 while instead we simulate these values based on the current posterior of $g(\cdot, \cdot)$. 

% Conditioning on these values when computing $\mu_*^{(n)}$ will make it so that retaining them in $\mathcal{S}$ provides no additional value, and will prevent us from doing this.

%The expected loss that would occur if we could observe $g(x,C(\S))$ is $\loss_n(x,C(\S))$.
%We then value observations at $x,\S$ according to the reduction in loss between this quantity 
%and the reduced loss $\loss_n(x,C(\S)\cup \S)$.

%Observe that 
%$\mathbb{E}_n[g(x',1) \vert g(x,C(\mathcal{S}))]$
%is the conditional expected value of the solution $x'$ given $g(x,\cdot)$ at fidelities in $C(\mathcal{S})$ under the time-$n$ posterior.  We then define $\mu_*^{(n,\emptyset)}$ to be the minimum over $x'$ of this quantity.
%This is the conditional expected loss associated with the best solution that could be obtained with this information.  We define $\mu_*^{(n + 1,\emptyset)} := \min_{x'} \mathbb{E}_{n+1}[g(x',1) \vert g(x,C(\mathcal{S}))]$ similarly.

%We then define a corresponding value of information and acquisition function,
%\begin{eqnarray}
%\label{eqn:takg0}
%\VOIE(x, \mathcal{S}) &=& \mathbb{E}_{n}\left[ \mu_*^{(n,\emptyset)} - \mu_*^{(n+1,\emptyset)} \vert x, \mathcal{S}\right]\nonumber\\
%\taKGE(x, \mathcal{S}) &=& \frac{\VOIE(x, \mathcal{S})}{\text{cost}^{(n)}(x , \max \mathcal{S})}
%\end{eqnarray}
%The expectation $\mathbb{E}_{n}$ is taken over both the values of $g(x,C(\mathcal{S}))$ (on which $\mu_*^{(n,\emptyset)}$ and $\mu_*^{(n+1,\emptyset)}$ depend), and also the value of $Y^{(n+1)}$ (on which only $\mu_*^{(n+1,\emptyset)}$ depends), both of which we simulate according to the time-$n$ posterior.

We define the $\taKGE$ acquisition function using this modified value of information as
\begin{equation}
\taKGE_n(x, \mathcal{S}) = \frac{\VOIE_n(x,\S)}{\text{cost}_n(x , \max \S)}
\label{eqn:takg0}
\end{equation}
To find the point $x$ and fidelity $s$ to sample, we optimize $\taKGE$ over 
% $x$ and $\S$ with fixed cardinality, and set $s$ to be the componentwise maximum of $\mathcal{S}$. 
$x$, fidelity $s$, and additional lower-fidelity point(s) $\S \setminus \{s\}$ as we did in \eqref{eqn:max-taKG}.

We refer to this VOI and acquisition function as ``0-avoiding,'' because they place 0 value on fidelities with any component equal to 0. 
% Indeed, for any fidelity $s$ with a component equal to $0$ and any collection of fidelities $\S \in B(s)$ resulting from observing $s$, we have $\S \subseteq C(\S)$ and $\taKGE_n(x,\S)=0$.
This prevents sampling at these points as long as the cost of sampling is strictly positive.

Indeed, suppose $s=\max(\S)$ has a component equal to $0$.  Then 
each element in $\S$ will have one component equal to $0$, and $\S \subseteq C(\S)$.
Then $\VOIE_n(x,\S) = \loss_n(x,C(\S)) - \loss_n(x,C(\S)\cup \S) = 0$.
Moreover, the following proposition shows that if $s=\max(\S)$ has all components strictly positive and additional regularity conditions hold, then $\VOIE_n(x,\S)$ is also strictly positive.

\begin{proposition}
If $s=\max(\S)$ has all components strictly positive, our kernel $K_n$ is positive definite, and the hypothesis of Proposition 1 is satisfied for $K_n$ given $g(x,C(\S))$, then $\VOIE_n(x,\S)$ is strictly positive.
\end{proposition}

Thus, maximizing $\taKGE$ will never choose to sample at a fidelity $s$ with a $0$ component.
Additionally, under other regularity conditions (see Corollary~1 in the supplement), $\VOIE_n(x,\S)$ is continuous in $\S$, and so the property that $\VOIE_n(x,\S)=0$ when a component of $s=\max(S)$ is 0 also discourages sampling at $s$ whose smallest component is {\it close} to 0.

%(In more detail, $\VOIE(x,\mathcal{S}) \ge 0$ because 
%\begin{align*}
%&\mathbb{E}_n \left[\mu_*^{(n,x,\mathcal{S})} - \mu_*^{(n+1,x,\mathcal{S})} \vert x, \mathcal{S}, g(x,C(\mathcal{S})) \right]\\
%&= \mu_*^{(n,x,\mathcal{S})} - \mathbb{E}_n \left[\mu_*^{(n+1,x,\mathcal{S})} \vert x, \mathcal{S}, g(x,C(\mathcal{S})) \right]
%\ge 0
%\end{align*}
%where the inequality is due to Jensen's inequality and the fact from iterated conditional expectation that 
%$\mathbb{E}_n[\mathbb{E}_{n+1}[g(x',1) \vert g(x,C(\mathcal{S}))] \vert g(x,C(\mathcal{S}))] = \mathbb{E}_n[g(x',1) \vert g(x,C(\mathcal{S}))]$.
%Taking the expectation of the last line in this equation block under the time-$n$ posterior then shows that $\VOIE(x, \mathcal{S}) \ge 0$.)

%\pfcomment{If we have time, might be good to put this paragraph in a proposition.}

%In a setting where the value of information becomes small as we approach $0$, this is an extremely useful aspect of the algorithm.  \pfcomment{I need to edit around here.}

%Although this completes theoretical specification of the $\taKG$ and $\taKGE$ acquisition function, 
%maximizing $\taKG$ and $\taKGE$ are challenging optimization problems and 
%naive brute-force approaches are unlikely to produce high-quality results.
%Thus, in the next section, we discuss efficient computational methods for solving these problems.


%Similar to Sect.~\ref{sect:mvoi}, to fix the issue that zero fidelity point provides positive value, we enhance the $\taKG$ acquisition function as
%\begin{eqnarray}
%\VOIE(x, s_{(1)}, \hat s_{(1)}, s_{(-1)}) &=& \mathbb{E}_{n}\left[\mu^{(n+2m-1)}_* \big\vert  
%\mathcal{E}\left(x, s_{(1)}, \hat s_{(1)}, s_{(-1)}\right) \right] \nonumber\\
%&& - \mathbb{E}_{n}\left[\mu^{(n+2m+1)}_* \Bigg\vert  
%\mathcal{E}\left(x, s_{(1)}, \hat s_{(1)}, s_{(-1)}\right) \cup \mathcal{S}(x, s_{(1)}, \hat s_{(1)}, s_{(-1)})
%\right]\nonumber\\
%\taKGE(x, s_{(1)},\hat s_{(1)}, s_{(-1)}) &=& \frac{\text{VOI+}(x, s_{(1)}, \hat s_{(1)}, s_{(-1)})}{\text{cost}^{(n)}(x ,s)},
%\label{eqn:taKG}
%%\end{split}
%\end{eqnarray}
%where $\mathcal{E}\left(x, s_{(1)}, \hat s_{(1)}, s_{(-1)}\right)$ is called the enhancing set and defined as $\mathcal{E}\left(x, s_{(1)}, \hat s_{(1)}, s_{(-1)}\right) := \langle x, s_{(1)}, s_{(-1)} \rangle \cup \langle x, \hat s_{(1)}, s_{(-1)}\rangle$ where
%\begin{eqnarray*}
%\langle x, s \rangle := \cup_{1 \le i \le m} \left\{(x, s^{-i})\right\}; s^{-i}_j = \left\{
%\begin{array}{ll}
%s_j & j \neq i\\
%0 & j = i.
%\end{array}
%\right.
%\end{eqnarray*}
%The subscription $j$ above means the $j$th dimension of the point. We visualize the enhancing set of a $2$d case in Fig.~\ref{Fig_enhancing_set}. The condition on the enhancing set is interpreted as we were to told the true value of the points in the set, while the condition on $\mathcal{S}(x, s_{(1)}, \hat s_{(1)}, s_{(-1)})$ is interpreted as we were to sample at the points there and observed the value with noise. The enhanced definition of $\VOIE$ quantifies the marginal benefit of evaluating at $\left(x, s_{(1)}, s_{(-1)}\right)$ and retaining $\left( x, \hat s_{(1)}, s_{(-1)}\right)$ after knowing the true values of some related zero fidelity points.
%

%Although this acquisition function considers the expected value of an improvement due to sampling, it differs from expected improvement approaches such as \citet{lam2015multifidelity} because the point at which an improvement occurs, $\argmax_{x\in\mathbb{A}} \mu^{(n+2)}(x,1_m)$ may differ from the point sampled.  Moreover, this acquisition function allows joint valuation of both the point $x$ and the fidelity $s$, while approaches such as \citet{lam2015multifidelity} require valuing a point $x$ assuming it will be evaluated at full fidelity and then choose the fidelity in a second stage.

%\subsection{Stochastic Gradients of $\loss_n(x,\S)$}
%\label{sect:estimate}

\subsection{Efficiently maximizing $\taKG$ and $\taKGE$}
\label{sect:optimize}

%We now describe simulation-based estimation of $\taKG$ and $\taKGE$.  Our
%efficient computational method described in the next section will levarage this
%particular estimation procedure to efficiently create unbiased gradient
%estimators of the acquisition function.

The $\taKG$ and $\taKGE$ acquisition functions are defined in terms of a hard-to-calculate function $\loss_n(x,\S)$. 
Here, we describe how to efficiently maximize these acquisition functions using stochastic gradient ascent with multiple restarts.  The heart of this method is a simulation-based procedure for simulating a stochastic gradient of $\loss_n(x,\S)$, i.e., a random variable whose expectation is the gradient of $\loss_n(x,\S)$ with respect to $x$ and the elements of $\S$.

To construct this procedure, we first provide a more explicit expression for $\loss_n(x,\S)$.
Because $\loss_n(x,\S)$ is the expectation of the minimum over $x'$ of $\E_n\left[g(x',1) \mid \Y(x,\S)\right]$, 
we begin with the distribution of this conditional expectation for a fixed $x'$ under the time-$n$ posterior distribution.

This conditional expectation can be calculated with GP regression from previous observations, the new point $x$ and fidelities $\S$, and the observations $\Y(x,\S)$.  This conditional expectation is linear in $\Y(x,\S)$.

Moreover, $\Y(x,\S)$ is the sum of $g(x,\S)$ (which is multivariate normal under the posterior) and optional observational noise (which is independent and normally distributed), and so is itself multivariate normal.
%under the time-$n$ posterior distribution.
As a multivariate normal random variable, it can be written as the sum of its mean vector and the product of the Cholesky decomposition of its covariance matrix with an independent standard normal random vector, call it $W$.  (The coefficients of this mean vector and covariance matrix may depend on $x$, $\S$, and previously observed data.)  The dimension of $W$ is the number of components in the observation, $|\S|$.

Thus, the conditional expected value of the objective $\E_n\left[g(x',1) \mid \Y(x,\S)\right]$ is a linear function (through GP regression) of another linear function (through the distribution of the observation) of $W$.

We also have that the mean of this conditional expectation is
$\E_n[\E_n[g(x’,1) | \Y(x,\S) ] ] =
\E_n[g(x’,1) ] = \mu_n(x’)$ by iterated conditional expectation.

These arguments imply the existence of a function $\tilde{\sigma}_n(x',x,\S)$ such that
$\E_n[g(x',1) | \Y(x,\S)]
= \mu_n + \tilde{\sigma}_n(x',x,\S) W$ simultaneously for all $x'$.
In the supplement, we show 
$\tilde{\sigma}_n(x',x,\S)= K_n\left((x', 1), x_{\S}\right) (D_{n}^T)^{-1}$
where $x_\mathcal{S} = \{(x, s): s \in \mathcal{S}\}$, and $D_{n}$ is the Cholesky factor of the covariance matrix $K_n\left(x_\mathcal{S}, x_\mathcal{S}\right)+\sigma^2 I$.
%$\mbox{diag}\left\{ \sigma_{i}^{2}\right\}$, and $\sigma_{i}^{2}$ is the variance of the noise of $g$ evaluated at the $i$th point in $x_\mathcal{S}$.

%where $\tilde{\sigma}_n(x’,x,\S)=\E_n[\mbox{Var}_{n+1}[g(x’,1)]|x^{(n+1)}, \mathcal{S}^{(n+1)}]$. \pfcomment{this definition of $\tilde{\sigma}$ is only correct in the noise-free case.}

Thus, 
\begin{equation*}
\loss_{n}\left(x,\S\right)=\mathbb{E}_{n}\left[\mbox{min}_{x'} \mu^{\left(n\right)}\left(x',1\right)+\tilde{\sigma}_{n}\left(x', x, \S\right)W\right].
\end{equation*}

% The notation $\diag{(0_{2m-1}, \sigma^2 1_{2})}$ is a diagonal matrix with the first $2m-1$ items zero and the last two items $\sigma^2$. The zero here indicates that we obtain the true value for points in the enhancing set and the $\sigma^2$ means that the two points we were to sample and retain are observed with noises.

When certain regularity conditions hold,
\begin{align*}
\nabla \loss_n(x,\S)
&= \nabla \mathbb{E}_{n}\left[\min_{x'}\left(\mu_n\left(x',1\right)+\tilde{\sigma}_{n}\left(x', x, \S\right)W\right)\right]\\
&= \mathbb{E}_{n}\left[\nabla \min_{x'}\left(\mu_n\left(x',1\right)+\tilde{\sigma}_{n}\left(x', x, \S\right)W\right)\right]\\
&= \mathbb{E}_{n}\left[\nabla \left(\mu_n\left(x^*,1\right)+\tilde{\sigma}_{n}\left(x^*, x, \S\right)W\right)\right]\\
&= \mathbb{E}_{n}\left[\nabla \tilde{\sigma}_{n}\left(x^*, x, \S\right)W\right],
\end{align*}
where $x^*$ is a global minimum (over $x'\in \mathbb{A}$) of $\mu_n(x',1) + \tilde{\sigma}_n(x',x,\S)W$, and the gradient in the last line is taken holding $x^*$ fixed even though in reality its value depends on $x$.
Here, the interchange of expectation and the gradient is justified using results from infinitessimal perturbation analysis \citep{l1990unified} and ignoring the dependence of $x^*$ on $x$ is justified using the envelope theorem \citep{milgrom2002envelope}.
We formalize this below in Theorem~\ref{stochastic_gradient}.

%Thus, to create an unbiased and strongly consistent estimator of the $\taKGE$ acquisition function, we rely on the Monte-carlo simulation, we sample a standard $|C(\mathcal{S})| + |\mathcal{S}|$-dimensional random vector $W$ to obtain our estimator $\hat{\loss}_{n}\left(x,Z_{(1:|C(\mathcal{S})|)}, W_{1:|C(\mathcal{S})|}\right) - \hat{\loss}_{n}\left(x,Z, W\right)$, where $\hat{\loss}_{n}\left(x,z, w\right)$ is defined by $\mbox{min}_{x'\in\mathbb{A}}\left(\mu^{\left(n\right)}\left(x',1_{m}\right)+\tilde{\sigma}_{n}\left(x',x,z\right)w\right)$, and we can repeat many times and take average,
%\begin{equation*}
%\hat{\loss}_{n}\left(x,z,\left\{w^j\right\}_j\right)=\frac{1}{k}\sum_{j=1}^{k} \mbox{min}_{x'\in\mathbb{A}}\left(\mu^{\left(n\right)}\left(x,1_{m}%%\right)+\tilde{\sigma}_{n}\left(x,x',z\right)w^{j}\right),
%\end{equation*}
%which is then a strongly consistent estimator by the law of large numbers.

% . By (\ref{eq:comp_qtaKG}) and the strong law of large numbers, we have that if $\left\{ W_{q}^{m}\right\} _{m\geq1}^{\infty}$ is a sequence of independent standard $q$-dimensional normal random vectors, then 

% \[
% \qtaKG (z_{1:q})=\underset{m\rightarrow\infty}{\mbox{lim}}\qtaKG_{m}(z_{1:q})\mbox{ a.s.},
% \]
% where
% \begin{equation}
% \begin{split}
% &\qtaKG_{m}(z_{1:q})\\
% &=\frac{\mu^{(n)}_* - \frac{1}{m}\sum_{j=1}^{m}\mbox{min}_{x\in\mathbb{A}}\left(\mu^{\left(n\right)}\left(x,1_{m}\right)+\tilde{\sigma}_{n}\left(x,z_{1:q}\right)W_{q}^{j}\right)}{\mbox{max}_{1\leq i\leq q}\mbox{cost}^{\left(n\right)}\left(x_{i},s_{i}\right)}.\\
% \end{split}
% \end{equation}
% Thus, $\qtaKG_{m}(z_{1:q})$ is an unbiased and strongly consistent estimator of $\qtaKG(z_{1:q})$.




%We will now describe efficient computation and maximization of the $\taKG$ and $\taKGE$ acquisition functions over subsets $\mathcal{S}$ of fixed cardinality, as specified in \eqn{max-taKG}.  We describe our method in detail only for $\taKGE$ as the method for $\taKG$ is similar. %We describe our method in the context of \eqn{max-taKG}. %since \eqn{max-q-taKG} follows the same ideas but the notation is much difficult to understand.
%Several challenges arise when solving these problems by adopting previously utilized methods.
%First, $\mathbb{A}$ is continuous, requiring optimization over a continuous domain when calculating term $\min_{x \in \mathbb{A}} \mu^{(n+q)}((x, 1))$ in Eq.~\eqn{q-taKG}, and then integration of this optimal value through its dependence on $y^{(n+1:n+q)}$.  Previous approaches (cite) have used discretization of $mathbb{A}$ to solve this problem, which creates 
%First, while \eqn{max-taKG} can be calculated analytically in terms of the normal cumulative distribution function (cdf), calculation of \eqn{max-q-taKG} for $q>1$ requires calculating multivariate normal cdfs of 
%To overcome these challenges, 

%We then use stochastic gradient ascent to optimize $\taKGE$ with multiple starts.

%We first introduce the following notation: $Z := C(\mathcal{S}) \cup \mathcal{S}$. 

\begin{theorem}
\label{stochastic_gradient}
Suppose $\mathbb{A}$ is compact, $\mu_0$ is constant, the kernel $K_0$ is continuously differentiable, and $\mathrm{argmin}_{x'\in\mathbb{A}}\ \mu_{n}(x', 1)+\tilde{\sigma}_n\left(x', x, \S\right)W$ contains a single element almost surely. Then, $\nabla \loss_n(x,\S) = \mathbb{E}_n\left[ \nabla \tilde{\sigma}(x^*,x,\S) W \right]$
%\[
%\nabla \loss_{n}(x,Z) =  \underset{r\rightarrow\infty}{\lim}\frac{1}{r}\sum_{j=1}^{r}\nabla\tilde{\sigma}_{n}\left(x^*_{j}, x, Z\right)W^{j}
%\]
%where $x^*_{j}$ is a global minimum in $\mathbb{A}$ of $\left(\mu_{n}(x', 1)+\tilde{\sigma}_n\left(x', x, Z\right)W^{j}\right)$ over $x'$.
%\begin{eqnarray*}
%F(z_{1:q})=\min_{x \in \mathbb{A}}\bmu{n}\left(x, 1\right)\\
%-\mathbb{E}_n\left[\min_{x \in \mathbb{A}} \left(\bmu{n} \left(x, 1\right)+\tilde{\sigma}_n\left(x, z_{1:q}\right)W_q\right)\right]
%\end{eqnarray*}
%and
\end{theorem}

With this result in place, we can obtain an unbiased estimator of $\nabla \loss_n(x,\S)$ by simulating $W$, calculating $x^*$, and then returning $\nabla \tilde{\sigma}_n(x^*,x,\S) W$.
Using this, together with the chain rule and an exact gradient calculation for $\text{cost}_n(x,\max \mathcal{S})$,
we can then compute stochastic gradients for $\taKG$ and $\taKGE$.

We then use this stochastic gradient estimator within stochastic gradient ascent \citep{kushner2003stochastic}
to solve the optimization problem \eqn{max-taKG} (or the equivalent problem for maximizing $\taKGE$).  The following theorem shows that, under the right conditions, a stochastic gradient ascent algorithm converges almost surely to a critical point of $\taKGE$. Its proof is in the supplement. 

\begin{theorem}
Assume the conditions of Theorem \ref{stochastic_gradient}, $\mathbb{A}$ is a compact hyperrectangle, and $\cost_n(\max \mathcal{S})$
is continuously differentiable and bounded below by a strictly positive constant. In
addition, assume that we optimize $\taKGE$ using a stochastic gradient
ascent method with the stochastic gradient from Theorem \ref{stochastic_gradient} whose stepsize
sequence $\left\{ \epsilon_{t}:t=0,1,\ldots\right\} $ 
satisfies $\epsilon_{t}\rightarrow0$, $\epsilon_{t}\geq0$,
$\sum_{t}\epsilon_{t}=\infty$ and $\sum_{t}\epsilon_{t}^{2}<\infty$.
Then the sequence of points $\left\{ x_{t}, \S_{t}\right\} _{t\geq0}$ from stochastic gradient ascent
converges almost surely to a connected set of stationary points of $\taKGE$.
\label{t:convergence}
\end{theorem}

\subsection{Warm-starting from partial runs}
\label{sect:warm-start}
When tuning hyperparameters using training iterations as a fidelity,
we can cache the state of training after $s$ iterations, for a warm start, and then continue training later up to a larger number of iterations $s'$ for less than $s'$ training iterations would cost at a new $x$.

We assume trace fidelities can be ``warm-started'' while non-trace fidelities cannot.  We also assume the incremental cost of evaluting fidelity vector $s'$ warm-starting from $s$ is the difference in the costs of their evaluations from a ``cold start''.  We model the cost of cold-start evaluation as in \sectn{sec:obj} with a Gaussian process on log(cost).  To obtain training data for this model, costs observed from warm-started evaluations are summed with those of the previous evaluations they continue to approximate what the cold-start cost would be.  We set $\cost_n(x,s)$ to be the difference in estimated cold-start costs if a previous evaluation would allow warm starting, and to the estimated cold start cost if not.

While our approach to choosing $x$ and $s$ to evaluate is to optimize $\taKGE$ as before 
our computational approach from \sectn{sect:optimize} required that $\cost_n(x,s)$ be continuously differentiable (in Theorem~\ref{t:convergence}).  This requirement is not met.
To address this, we modify the way we optimize the $\taKGE$ acquisition function.
(The approach we describe also works for $\taKG$.)

First, we maintain a basket of size at most $b$ of previously evaluated point, fidelity pairs, $(x(j),s(j))$.
For each $j\le b$, we optimize $\taKGE_n(x(j),\S)$ letting $\S$ vary over those sets satisfying two conditions: (1) $|\S| = \ell$; 
(2) $s' \ge s(j)$ componentwise for each $s' \in \S$, 
with equality for non-trace fidelity components. 
%(2) the non-trace fidelities of each $s' \in \S$ match those of $s(x)$; (3) the trace fidelities of each $s' \in \S$ are bounded below by those of $s(x)$.
Over this set, $\cost_n(x(j),\S)$ is continuously differentiable in $\S$ and the method from \sectn{sect:optimize} can be applied.

% This gives us an $\S^*(x)$ and a value $\taKG_n(x,\S^*(x))$.

We also optimize $\taKGE_n(x,\S)$ over all $x$ and $\S$ with $|\S|=\ell$, but using the estimated cold-start cost function and the method from \sectn{sect:optimize}.

Among the solution to these at most $b+1$ optimization problems, we select the $x$ and $\S$ that provide the largest $\taKGE_n(x,\S)$ at optimality, and evaluate $g$ at $x$ and $\max(\S)$.

We then update our basket.
We first add the $x$ and $\max(\S)$ produced by the optimization not constraining $x$.
If the basket size exceeds $b$, we then remove the $x$ and $s$ whose optimization over $\taKGE_n$ produced the smallest value.
In practice, we set $b=10$.

% To address this, we suppose without loss of generality that the fidelities are ordered so that the non-trace fidelities of $s$ are $s_{1:m_2}$.


%First, we keep a basket of size $b$ containing previously evaluated points $x$ and corresponding non-trace fidelities $s_{1:m_2}$ at which they were evaluated.  For the $j$th pair of $x$ and $s_{1:m_2}$ in the basket, we set $x(j)$ to $x$ and construct a fidelity vector $s(j)$.  The non-trace fidelities in $s(j)$ are simply equal to the corresponding value in $s_{1:m_2}$.  Each trace fidelity is equal to its largest previously evaluated value, among evaluations matching $x$ and $s_{1:m_2}$.  If $x$ has been evaluated at multiple non-trace fidelities, then each may appear separately in the basket.

%For each $j$, we identify a correspending restricted set of fidelity sets, 
%$\{ \S : |\S| = \ell, s \ge s(j), s_{1:m_2} = s_{1:m_2}(j)\ \forall s \in \S \}$.
%Over this set, $\cost_n(x(j),\S)$ is continuously differentiable and the method from \sectn{sect:optimize} can be applied.

% $x(j)$ and $s(j)$ in the basket, we maximize $\taKGE(x,\S)$, holding $x$ fixed, over those $\S$ whose elements are all componentwise larger than $s(j)$ and that satisfy the cardinality constraint $|\S|=\ell$.

%(Note that the restriction to $\S$ with elements componentwise larger than $s_n(x)$ is not without loss of generality.  If we were able to retain the full trace then we could drop this constraint and obtain the same solution, because evaluating a lower fidelity would have no value.

% When we finish the evaluation at iteration $n-1$, we first re-train the GP and re-compute the  posterior distribution. Now we compute the warm-start $\taKGE$ ($\text{WS-\taKGE}$) if we warm start from $(x, \max \mathcal{S})$ to $(x, \tilde{\mathcal{S}})$ by
% \begin{eqnarray*}
% &&\text{WS-\taKGE}(x, \tilde{\mathcal{S}})=\\
% &&\max_{\substack{\max\tilde{\mathcal{S}}^{m_1} \ge \max\mathcal{S}^{m_1} \\ \max\tilde{\mathcal{S}}^{m_2} = \max\mathcal{S}^{m_2}}} \frac{\VOIE(x, \tilde{\mathcal{S}})}{\text{cost}^{(n)}(x , \max \tilde{\mathcal{S}}-\max \mathcal{S})}
% \end{eqnarray*}
% where $\mathcal{S}^{m_1}$ and $\tilde{\mathcal{S}}^{m_1}$ are the trace fidelities in $\mathcal{S}$ and $\tilde{\mathcal{S}}$, and $\mathcal{S}^{m_2}$ and $\tilde{\mathcal{S}}^{m_2}$ are the non-trace fidelities. We can solve $b+1$ small-scale continuous optimization problems to obtain $\text{WS-\taKGE}(x, \tilde{\mathcal{S}})$ for all the points in the basket and the sample at the last iteration. We pick the largest $\text{WS-\taKGE}(x, \tilde{\mathcal{S}})$ as the potential warm-starting point and then drop the smallest one from $b+1$ candidates to obtain the new basket.

% Next we decide whether to warm start with the optimal $\text{WS-\taKGE}(x, \tilde{\mathcal{S}})$ or start a new point from scratch. This can be achieved by comparing $\max \text{WS-\taKGE}(x, \tilde{\mathcal{S}})$ and the maximum $\taKGE$, where $\taKGE$ is defined in \eqn{takg0}.

\subsection{Batch and Derivative Evaluations}
\label{sect:extensions}
$\taKG$ and $\taKGE$ generalize naturally 
following \citet{wu2016parallel} and \citet{wu2017bayesian}
to batch settings where we can evaluate multiple point, fidelity pairs simultaneously and derivative-enabled settings where we observe gradients. 

The batch version uses the same acquisition functions $\taKG$ and $\taKGE$ defined above, but optimizes over a {\it set} of values for $s$, each of which has an associated $\S \in B(s)$ of limited cardinality.

In the derivative-enabled setting, we incorporate (optionally noisy) gradient observations into our posterior distribution directly through GP regression.  We also generalize the $\taKG$ and $\taKGE$ acquisition functions to allow inclusion of gradients of the objective in the set of quantities observed $\Y(x,\S)$ in the definition of $L_n(x,\S)$.


\section{NUMERICAL EXPERIMENTS}
\label{sect:numerical}
We compare sequential, batch, and derivative-enabled $\taKGE$ with benchmark algorithms on synthetic optimization problems (Sect.~\sect{synthetic}), hyperparameter optimization of neural networks (Sect.~\sect{hyperexp}), and hyperparameter optimization for large-scale kernel learning (Sect.~\sect{kernel-learning}).  The synthetic and neural network benchmarks use fidelities with trace observations, while the large-scale kernel learning benchmark does not.  We integrate over GP hyperparameters by sampling $10$ sets of values using the \texttt{emcee} package \citep{foreman2013emcee}.


\subsection{Optimizing synthetic functions} \label{sect:synthetic}
Here, we compare $\taKGE$ against both the sequential and batch versions of the single-fidelity algorithms KG \citep{wu2016parallel} and EI \citep{jones1998efficient,wang2015parallel}, a derivative-enabled single-fidelity version of KG \citep{wu2017bayesian}, Hyperband \citep{li2016hyperband}, and the multi-fidelity method BOCA \citep{kandasamy2017multi}.  BOCA is the only previous method of which we are aware that treats multiple continuous fidelities.  We do not compare against FaBOLAS \citep{klein2016fast,klein2015towards} because the kernel it uses is specialized to neural network hyperparameter tuning.

Following experiments in \citet{kandasamy2017multi}, we augment four synthetic test functions, 2-d Branin, 3-d Rosenbrock, 3-d Hartmann, and 6-d Hartmann, by adding one or two fidelity controls, as described in the supplement.  
% We chose to include only a single additional fidelity control because adding more than one provided an even more substantial benefit of the multi-fidelity methods over single-fidelity ones.
%We use cost $\text{cost}(x, s) := \prod_{i=1}^{m} (0.01 + s_i)$ 
We set the cost of an individual evaluation of $x$ at fidelity $s$ to $0.01 + \prod_i s_i$.
Thinking of $s_1$ as mimicking training data and $s_2$ training iterations, 
the term $\prod_i s_i$ mimics a cost of training that is proportional to the number of times a datapoint is visited in training.  The term $0.01$ mimics a fixed cost associated with validating a trained model.
We set the cost of a batch evaluation to the maximum of the costs of the individual evaluations,
to mimic wall-clock time for running evaluations synchronously in parallel. 

%\newcommand{\figwidth}{95pt}
%\newcommand{\figheight}{100pt}
%\newcommand{\figsize}{95pt}
%\newcommand{\figwidth}{140pt}
%\newcommand{\figheight}{128pt}
%\newcommand{\figwidth}{180pt}
%\newcommand{\figheight}{140pt}
\newcommand{\figwidth}{230pt}
\newcommand{\figheight}{190pt}




\begin{figure*}[tb]
\centering
  \subfigure
  \centering
  \includegraphics[width=\figwidth, height = \figheight]{Branin.pdf}
  \subfigure
  \centering
  \includegraphics[width=\figwidth, height = \figheight]{Rosenbrock.pdf}\\
    \subfigure
  \centering
\includegraphics[width=\figwidth, height = \figheight]{Hartmann3.pdf}
   \subfigure
  \centering
  \includegraphics[width=\figwidth, height = \figheight]{Hartmann6.pdf}
\vspace{-8pt}
\caption{\small Optimizing synthetic functions: Plots show simple regret over $40$ independent runs for synthetic functions with trace observations and one or two continuous fidelity controls for 2-d Branin, 3-d Rosenbrock, 3-d Hartmann, and 6-d Hartmann problems.  $\taKG0$ performs well compared with a variety of competitors in both sequential and batch settings.}
\vspace{-12pt}
\label{Fig_syn_taKG}
\end{figure*}



Fig.~\ref{Fig_syn_taKG} shows results.  For methods that have both sequential and batch versions, the batch size is indicated with a number before the method name.  For example, 8-EI indicates expected improvement performed with batches of size 8.  We run versions of $\taKGE$ using $|\S|$ set to $2$ (taKG\_0 2-points) and $3$ (taKG\_0 3-points).

We first see that using the larger $|\S|$ improves the performance of $\taKGE$.
We then see that, for both values of $|\S|$, sequential $\taKGE$ performs well relative to sequential competitors (1-EI, 1-KG, 1-BOCA), and batch $\taKGE$ with batch size 8 (8-$\taKGE$) performs well relative to its batch competitors (8-EI, 8-KG, Hyperband).

Here, we consider Hyperband to be a batch method although the amount of parallelism it can leverage varies through the course of its operation.

\subsection{Optimizing hyperparameters of neural networks } \label{sect:hyperexp}
Here, we evaluate on hyperparameter optimization of neural networks. Benchmarks include the single-fidelity Bayesian optimization algorithms KG \citep{wu2016parallel} and EI \citep{jones1998efficient}, their batch versions, and the state-of-art hyperparameter tuning algorithms HyperBand and FaBOLAS.  $\taKGE$ uses two fidelity controls: the size of the training set and the number of training iterations.  

Following \citet{li2016hyperband}, we set the cost to the number of training examples passed during training divided by the number passed in full fidelity. For example, if we have $5\times 10^4$ training points and the maximum number of epochs is $100$, then the cost of evaluating a set of hyperparameters using $10^4$ sub-sampled training points per epoch over $10$ epochs is $10^4 \times 10 / (5\times10^4 \times 100) = 0.02$. Complete training has cost $1$.

\begin{figure*}[tb]
  \centering
  \subfigure
  \centering
  \includegraphics[width=\figwidth, height = \figheight]{MNIST.pdf}
  \subfigure
  \centering
  \includegraphics[width=\figwidth, height = \figheight]{CIFAR10.pdf}\\
  \subfigure
  \centering
  \includegraphics[width=\figwidth, height = \figheight]{SVHN.pdf}
  \subfigure
  \centering
  \includegraphics[width=\figwidth, height = \figheight]{KISSGP.pdf}
\vspace{-8pt}
\caption{\small We show the -log marginal likelihood divided by the number of datapoints for tuning feedforward neural networks on MNIST (each with 20 runs); tuning convolutional neural networks on CIFAR-10 and SVHN (each with 10 runs); and for KISS-GP kernel learning.
$\taKGE$ outperforms competitors in both sequential and batch settings.  }
\vspace{-12pt}
\label{Fig_deep_taKG}
\end{figure*}


\paragraph{Feedforward neural networks on MNIST}
\label{sect:mlps}
We tune a fully connected two-layer neural network on MNIST. The maximum number of epochs allowed is $20$. We optimize $5$ hyperparameters: learning rate, dropout rate, batch size and the number of units at each layer. 

Fig.~\ref{Fig_deep_taKG} shows that sequential $\taKGE$ performs much better than the sequential methods KG, EI and the multi-fidelity hyperparameter optimization algorithm FaBOLAS. $\taKGE$ with a batch size $4$ substantially improves over batch versions of KG and EI, and also over the batch method Hyperband.

\paragraph{Convolutional neural networks on CIFAR-10 and SVHN}
\label{sect:cnns}
We tune convolution neural networks (CNNs) on CIFAR-10 and SVHN. Our CNN consists of 3 convolutional blocks and a softmax classification layer. Each convolutional block consists of two convolutional layers with the same number of filters followed by a max-pooling layer. There is no dropout or batch-normalization layer. We split the CIFAR-10 dataset into 40000 training samples, 10000 validation samples and 10000 test samples. We split the SVHN training dataset into 67235 training samples and 6000 validation samples, and use the standard 26032 test samples. We apply standard data augmentation: horizontal and vertical shifts, and horizontal flips. We optimize $5$ hyperparameters to minimize the classification error on the validation set: the learning rate, batch size, and number of filters in each convolutional block. Hyperband uses the size of the training set as its resource (it can use only one resource or fidelity), using a bracket size of $s_{\max} = 4$ as in \citet{li2016hyperband} and the maximum resource allowed by a single configuration set to 40000.
We set the maximum number of training epochs for all algorithms to $50$ for CIFAR-10 and $40$ for SVHN. 
Because of the computational expense of training CNNs, we leave out some benchmarks, dropping the single-fidelity method EI in favor of the structurally similar single-fidelity method KG, and performing batch evaluations for only some methods.

Fig.~\ref{Fig_deep_taKG} shows that sequential $\taKGE$ outperforms its competitors (including the batch method Hyperband) on both problems.  Using batch evaluations with $\taKGE$ on CIFAR-10 improves performance even further.

When we train using optimized hyperparameters on the full training dataset for $200$ epochs, test data classification error is $\sim12\%$ for CIFAR-10 and $\sim5\%$ for SVHN. 


\subsection{Optimizing hyperparameters for large-scale kernel learning}
\label{sect:kernel-learning}
We test derivative-enabled $\taKGE$ (ta-dKG$^\emptyset$) in a large-scale kernel learning example: the 1-d demo example for KISS-GP \citep{wilson2015kernel} on the GPML website \citep{gpml}.  In this example, we optimize 3 hyperparameters (marginal variance, length scale, and variance of the noise) of a GP with an RBF kernel on $1$ million training points to maximize the log marginal likelihood. We evaluate both the log marginal likelihood and its gradient using the KISS-GP framework. We use two continuous fidelity controls: the number of training points and the number of inducing points.  We set the maximum number of inducing points to $m=1000$.

We compare ta-d-KG to the derivative-enabled knowledge gradient ($\qKGg$) \citep{wu2017bayesian}, using both algorithms in the sequential setting (1-dKG and 1-cf-dKG) and with a batch size of 4 (4-dKG and 4-cf-dKG).  We leave out methods that are unable to utilize derivatives, as these are likely to substantially underperform.

Fig.~\ref{Fig_deep_taKG} shows that ta-dKG$^\emptyset$ successfully utilizes inexpensive function and gradient evaluations to find a good solution more quickly than d-KG, in both the sequential and batch setting.

%which dramatically improves the default hyperparameters set in the tutorial ($14\%$). The details of the optimized hyperparameters and the default hyperparameters are included in the appendix.


% used to be 125pt by 132 pt --- making a little smaller to save space
% If we have more room, can do 130x135

\section{CONCLUSION}
\label{sect:conclusion}
We propose a novel multi-fidelity acquisition function, the trace aware knowledge-gradient, which 
leverages special structure provided by trace observations,
is able to handle multiple simultaneous continuous fidelities, and 
generalizes naturally to batch and derivative settings.
This acquisition function uses traces to find 
good solutions to global optimization problems more quickly than state-of-the-art algorithms
in application settings including deep learning and kernel learning.

% $\taKG$ and $\taKGE$, which generalize naturally to batch and derivative settings. $\taKGE$ 
%uses traces to
%find good hyperparameters for deep and kernel learning
%more quickly than state-of-art algorithms. 


%\subsubsection*{References}

\bibliographystyle{apalike}
\bibliography{pKG}

\newpage

%References follow the acknowledgements.  Use unnumbered third level
%heading for the references title.  Any choice of citation style is
%acceptable as long as you are consistent.

\section{SUPPLEMENTARY MATERIAL}

\subsection{Background: Gaussian processes}
\label{sect:gp}


We put a Gaussian process (GP) prior \citep{rasmussen2006gaussian} on the function $g$.  
%We describe this procedure placing the prior on $g$ directly, and then discuss below when we recommend instead placing it on $(x,s) \mapsto \log g(x,s)$.
The GP prior is defined by its mean function $\mu_{0}: \mathbb{A} \times [0, 1]^m \mapsto \mathbb{R}$ and kernel function $K_{0}: \left\{\mathbb{A} \times [0, 1]^m\right\} \times \left\{\mathbb{A} \times [0, 1]^m \right\} \mapsto \mathbb{R}$.  These mean and kernel functions have hyperparameters, whose inference we discuss below.

We assume that evaluations of $g(x,s)$ are subject to additive independent normally distributed noise with common variance $\sigma^2$.
% \begin{eqnarray*}
% h(x, s) | g(x, s) \sim \mathcal{N}(g(x, s), \sigma^2),
% \end{eqnarray*}
% where $h(x, s)$ is the value observed when we sample at $x$ with fidelity $s$. 
We treat the parameter $\sigma^2$ as a hyperparameter of our model, and also discuss its inference below.  Our assumption of normally distributed noise with constant variance is common in the BO literature \citep{klein2016fast}. 

The posterior distribution on $g$ after observing $n$ function values at points $z_{(1:n)} := \{(x_{(1)}, s_{(1)}), (x_{(2)}, s_{(2)}), \cdots, (x_{(n)}, s_{(n)})\}$ with observed values $y_{(1:n)} := \{y_{(1)}, y_{(2)}, \cdots, y_{(n)}\}$ remains a Gaussian process \citep{rasmussen2006gaussian}, and $g \mid z_{(1:n)}, y_{(1:n)} \sim \text{GP}(\mu_n, K_n)$ with $\mu_n$ and $K_n$ evaluated at a point $z=(x,s)$ (or pair of points $z$, $\tilde{z} = (\tilde{x},\tilde{s})$) given as follows

\begin{equation}
\begin{split}
& \bm\mu_n(\bm{x}) = \mu(\bm{x})\\
&\qquad + K(\bm{x},x_{1:n}) \left(K(x_{1:n},x_{1:n})+\right.\\
&\qquad \left.\sigma^2 I \right)^{-1}  (y_{1:n}-\mu(x_{1:n})),\\
& K_{n}(\bm{x_1}, \bm{x_2}) =  K(\bm{x_1},\bm{x_2})\\
&\qquad - K(\bm{x_1},x_{1:n})  \left(K(x_{1:n},x_{1:n})\right.\\
&\qquad\left.+\sigma^2 I\right)^{-1} K(x_{1:n},\bm{x_2}).
\end{split}
\label{eq:posterior}
\end{equation}

We emphasize that a single evaluation of $g$ provides multiple observaions of $g$, because of trace observations.  taKG chooses to retain 2 observations per evaluation, and so $n$ will be twice the number of evaluations.

%When $g$ is the validation error in a hyperparameter optimization problem, we recommend putting a GP prior on $\log g(x,s)$, rather than on $g(x,s)$ directly, because (1) $g(x,s)$ is nonnegative and will be allowed to be negative after log scaling, better matching the range of values assumed by the GP, and (2) because $g(x,s)$ can climb steeply over several orders of magnitude as we move away from the optimal $x$, making $\log g(x,s)$ easier to model. 

This statistical approach contains several hyperparameters: the variance $\sigma^2$, and any parameters in the mean and kernel functions.  We treat these hyperparameters in a Bayesian way as proposed in \citet{snoek2012practical}. We analogously train a separate GP on the logarithm of the cost of evaluating $g(x,s)$. 

%\section{Extension to batch settings}
%The taKG generalizes naturally to batch settings where we can evaluate multiple (point, fidelity) pairs at once. We value joint evaluation of $q \ge 1$ points $x_{1:q}$ at fidelities $s_{1:q}$ and retain $q$ middle points $\hat s_{1:q}$, where $z_{1:q} = ( (x_{1}, \hat s_{1}, s_{1}), \ldots, (x_q, \hat s_{q}, s_q))$ are the points we are proposing to evaluate and retain, by
%% pairs $z^{(n+1:n+q)}$, where $z^{(n+i)} = (x^{(n+i)}, s^{(n+i)})$, 
%\begin{equation}
%%\qtaKG (z^{(n+1:n+q)}) = \frac{\min_{x' \in \mathbb{A}} \mu^{(n)}(x', 1) - \min_{x' \in \mathbb{A}} \mu^{(n+q)}(x', 1) \mid z^{(n+1:n+q)}  \right]}{\max_{1 \le i \le q}\text{cost}^{(n)}(z^{(n+i)})},
%\qtaKGE(z_{1:q}) =\frac{\mathbb{E}_{n}\left[\mu^{(n+q(2m-1))}_* \mid \langle z_{1:q} \rangle \right] - \mathbb{E}_{n}\left[\mu^{(n+q(2m+1))}_* \mid \langle z_{1:q} \rangle, z_{1:q}  \right]}{\max_{1 \le i \le q}\text{cost}^{(n)}(x_{i},s_{i})},\nonumber\\
%\label{eqn:q-taKG}
%\end{equation}
%We then modify \eqn{max-taKG} by sampling at the batch of points and fidelities that maximize
%\begin{eqnarray}
%\max_{z_{1:q}} \qtaKGE (z_{1:q})
%\label{eqn:max-q-taKG}
%\end{eqnarray}
%
%\section{Extension to derivative-enabled settings}
%\label{sect:derivative}
%
%% When we can access the gradient information of the expensive function, BO methods can naturally incorporate this additional information to speed up the optimization process, as shown in \citep{wu2017bayesian}. However, the function and its gradient evaluation are usually both expensive. For example, in the large-scale Gaussian process kernel learning, one can compute the log likelihood and its gradient although the computation scales cubically with the number of data points; fortunately, one can have {\bf biased} estimations of both the log likelihood and its gradient by evaluating on a subset of the training data. Our method extends to settings where one can access both the function value and the gradient information.
%
%Following \citet{wu2017bayesian}, which developed single-fidelity BO methods for use when gradient information is available, we develop a version of the taKG algorithm, called derivative-enabled taKG, that can be used when gradients are available in the continuous-fidelity setting.  This can be used to reduce the number of function evaluations required to find a high-quality solution to \eqn{min_f}.  
%
%First, observe that a $GP(\mu^n, K^n)$ prior on $g(x, s)$ implies a multi-output $GP(\tilde \mu^n, \tilde K^n)$ prior on $(g, \nabla_x g)$. Then, observe that the same reasoning we used to develop the taKG acquistion function in \eqn{taKG} can be used when when we observe gradients to motivate the acquisition function,
%\begin{equation*}
%\begin{split}
%\text{ta-d-KG} (x, s) 
%= \frac{\tilde \mu_{*}^{(n)} - \mathbb{E}_{n}\left[\tilde \mu_{*}^{(n+2)}\bigg\vert 
%x, 
%s_{(1)},
%\hat s_{(1)}, s_{(-1)}
%\right]	}{\text{cost}^{(n)}(x ,s)},
%\end{split}
%\end{equation*}
%where $\tilde \mu_{*}^{(n)}$ and $\tilde \mu_{*}^{(n+2)}$ are the optimal mean predictions, $\text{cost}^{(n)}(x, s)$ is now our estimate after $n$ samples of the cost of evaluating both $g$ and its gradient with respect to $x$,
%and $\tilde{\mu}_1^{(n)}(x,s)$ is the posterior mean on $g(x,s)$ in the multi-output GP. 
%The conditional expectation is taken with respect to both the observed function value and gradient.
%A batch version of the derivative-enabled taKG acquistion function can be defined analogously. We use techniques similar to those described in Sect.~\ref{sect:compute} to optimize these acquisition functions.

% \citet{kandasamy2016multi} generalizes the upper confidence bound (UCB) criteria to the discrete-fidelity setting, and further generalizes to continuous fidelities in \citet{kandasamy2017multi}.  We compare against \citet{kandasamy2017multi} in our numerical experiments, finding that taKG provides improved performance when the budget is sufficiently large.  This may be because the UCB criteria was originally designed for minimizing the cumulative regret and not simple regret, with the simple regret bounds resulting from theoretical analysis of cumulative regret being loose enough to be consistent with the empirical performance gap we observe in our experiments.  It may also be because both \citet{kandasamy2016multi} and \citet{kandasamy2017multi} use a two-stage process in which the point to evaluate is selected without considering the fidelity, and the fidelity only selecting afterward.  In contrast, taKG selects the point and fidelity jointly, and considers the impact of fidelity choice on the best point to sample.
%
% \citet{poloczek2016multi} develops a knowledge gradient method for the discrete-fidelity setting, but does not consider the continuous-fidelity setting.  Our computational techniques are quite different from the ones developed there, as necessitated by our consideration of continuous fidelities.  As an advantageous byproduct, our computational techniques also avoid the need to discrete $\mathbb{A}$.
%
% This paper proposes the first knowledge gradient method for continuous-fidelity settings, and show how to generalize it to the batch and derivative-enabled settings by generalizing the computational technique developed in \citet{wu2017bayesian}. To the best of our knowledge, this is the first multi-fidelity batch Bayesian optimization algorithm. 

%\section{Warm-starting from partial runs}
%When choosing the point for the next iteration, we sometimes would like to warm start from where left in the previous iterations if we can observe the trace. However, evaluating further at some previously evaluated point will cause the discontinuity of the cost function at these points, and the computational method in Sect.~\sect{compute} will fail. In this section, we will propose a computationally efficient and principled way to choose whether warm-starting from a previous point or starting a new point from scratch.
%
%We first focus on the fully-sequential setting where we will propose a single point at each iteration. We keep a basket of $b$ evaluated points, practically one can choose $b=10$. When we finish the evaluation at iteration $n-1$, we first re-train the GP and re-compute the  posterior distribution. Now we compute the maximum \emph{Value of Information} (VOI) per unit cost if we warm-start from either any point in the basket or the point just evaluated at iteration $n-1$ by
%\begin{equation*}
%\begin{split}
%\text{Warm-Start}(x,&s_{(-1)}) = \max_{s_{(1)}, \hat s_{(1)} \le s_{(1)}} \taKG (x, s_{(1)}, \hat s_{(1)}, s_{(-1)}) 
%\end{split}
%\end{equation*}
%where $x$ and $s_{(-1)}$ are evaluated points and the corresponding other fidelity controls. We obtain $\text{Warm-Start}(x, s_{(-1)})$ by solving $b+1$ small-scale continuous optimization problems .We pick the largest $\text{Warm-Start}(x, s_{(-1)})$ as the potential warm-starting point if we choose to by solving and drop the smallest one from $b+1$ candidates to obtain the new basket.
%
%Next we need to decide whether to warm start using  $\text{Warm-Start}(x, s_{(-1)})$ or start a new point from scratch. This can be achieved by comparing $\max \text{Warm-Start}(x, s_{(-1)})$ and the maximum VOI of starting a new point from scratch. The maximum VOI of starting a new point from scratch is defined as $\max_{x, s_{(1)}, \hat s_{(1)}, s_{(-1)}} \taKG (x, s_{(1)}, \hat s_{(1)}, s_{(-1)})$ where $\taKG (x, s_{(1)}, \hat s_{(1)}, s_{(-1)})$ is defined in \eqn{taKG} and can be solved by the computational method in Sect.~\sect{compute}.
%
%Now we turn to batch settings. We select the first point by the method described above, and denote the solution as $(x^{*}, s^{*}_{(1)}, \hat s^{*}_{(1)}, s^{*}_{(-1)})$ where $(x^{*}, s^{*}_{(-1)})$ is either an evaluated point or a new point and $s^{*}_{(1)}$ is the proposed new training length and $\tilde s^{*}_{(1)}$ is the proposed new middle point to retain. For the next $q-1$ points, we choose to all start from scratch. To get these points, we will fix the first point $(x^{*}, s^{*}_{(1)}, \hat s^{*}_{(1)}, s^{*}_{(-1)})$, and then select the next $q-1$ points by optimizing $\max_{\z{1:(q-1)}} \taKG(\z{1:(q-1)}, (x^{*}, s^{*}_{(1)}, \hat s^{*}_{(1)}, s^{*}_{(-1)}))$ by the method in Sect.~\sect{compute}.

\subsection{Proofs Details}

In this section we prove the theorems of the paper. We may assume without loss of generality that $|\mathcal{S}| = 1$, and we denote the number of fidelities by $m$. Observe that the dimension of the vector  $C(\mathcal{S}) \cup \mathcal{S}$ is $q:=2m$.
We first show some smoothness properties of $\tilde{\sigma}_{n}$, $\mu_{n}$ and  $\mbox{cost}_{n}$ in the following lemma.
\begin{lemma}
\label{smooth_lemma}
We assume that the domain  $\mathbb{A}$ is compact, $\mu_0$ is a constant and the kernel $K_{0}$ is continuously differentiable. We then have that
\begin{enumerate}
\item $\mu_{n}$, and $\tilde{\sigma}_{n}\left(\cdot,z_{1:q}\right)$ are both continuously differentiable for any vector $z_{1:q}$.
\item For any $x'$, $\tilde{\sigma}_{n}\left(x',z_{1:q}\right)$ is continuously  differentiable respect to $z_{1:q}$.
\item $\mbox{cost}_{n}$ is continuously differentiable.
\item $\max_{1 \le i \le q}\text{cost}_{n}(x_i, s_i)$ is differentiable if $\left|\mbox{arg max}_{1\leq i\leq q}\mbox{cost}_{n}\left(x_{i},s_{i}\right)\right|=1$.
\end{enumerate}
\end{lemma}
\proof{}
The posterior parameters of the Gaussian process $\mu_{n}$, $K_{n}$ and $\mbox{cost}_{n}$ are continuously differentiable if the kernel $K_{0}$ and $\mu_{0}$ are both continuously differentiable. Observe that $\tilde{\sigma}_n(x, z_{(1:q)})= K_{n}\left((x, 1), z_{1:q}\right) (D_{n}\left(z_{1:q}\right)^T)^{-1}$, and so $\tilde{\sigma}_{n}\left(\cdot,z_{1:q}\right)$ is continuously differentiable because $K_{n}$ is continuously differentiable. This proves (1) and (3). (4) follows easily from (3).

To prove (2) we only need to show that $(D_{n}\left(z_{1:q}\right)^T)^{-1}$ is continuously differentiable respect to $z_{1:q}$. This follows from the fact that multiplication, matrix inversion (when the inverse exists), and Cholesky factorization \citep{smith1995differentiation} preserve continuous differentiability. This ends the proof.
\endproof


We now prove Theorem 1.
\proof{}
Let $W_q$ be a standard normal random vector, and define $f(x,z_{1:q}):=\mu_n \left(x, 1\right)+\tilde{\sigma}_n\left(x, z_{1:q}\right)W_q$. By Lemma \ref{smooth_lemma}, $f$ is continuously differentiable. By the envelope theorem (see Corollary 4 of \citealt{milgrom2002envelope}), $\nabla f(y,z_{1:q})=\nabla \tilde{\sigma}_n\left(y, z_{1:q}\right)W_q$ a.s., where $y=\mbox{arg max}_{x\in \mathbb{A}}\left(\mu_n(x, 1_q)+\tilde{\sigma}_n\left(x, z_{1:q}\right)W_{q}\right)$.

We now show that we can interchange the gradient and the expectation in $\nabla \mathbb{E}_n\left[\min_{x \in \mathbb{A}} \left(\mu_{n} \left(x, 1\right)+\tilde{\sigma}_n\left(x, z_{1:q}\right)W_q\right)\right]$. Observe that the domain of $z_{1:q}$ is compact,  $\tilde{\sigma}_{n}\left(x',z_{1:q}\right)$ is continuously  differentiable respect to $z_{1:q}$ by Lemma \ref{smooth_lemma}, and so $\left\Vert\tilde{\sigma}_{n}\left(x',z_{1:q}\right)\right\Vert$ is bounded. Consequently, Corollary 5.9 of \citet{bartle} implies that we can interchange the gradient and the expectation. The statement of the theorem follows from the strong law of large numbers.
\endproof

The following corollary follows from the previous proof.

\begin{corollary}
Under the assumptions of the previous theorem, $L_{n}(x,S)$ is continuous.
\end{corollary}

We now prove Theorem 2.
\proof{}
We prove this theorem using Theorem 2.3 of Section 5 of \citet{kushner2003stochastic}, which depends on the structure of the stochastic gradient  $G$ of the objective function. In addition, we simplify the notation and denote $(X,\S)$  by $Z$.

% Thus, first we observe that from Theorem 1, the expression of the stochastic gradient of $taKG+$ is:
% \begin{eqnarray*}
% G\left(Z\right) & =\\
% \frac{1}{r}\sum_{k=1}^{r}\frac{C\left(Z\right)\nabla\tilde{\sigma}_{n}\left(y_{k},Z\right)W_{q}^{k}}{C\left(Z\right)^{2}}\\
% -\frac{1}{r}\sum_{k=1}^{r}\frac{\nabla C\left(Z\right)}{C\left(Z\right)^{2}}\left(\mu^{\left(n\right)}\left(y_{k},1_{m}\right)+\tilde{\sigma}_{n}\left(y_{k},Z\right)W_{q}^{k}\right)
% \end{eqnarray*}
% where $C\left(Z\right):=\mbox{max}_{1\leq i\leq q}\mbox{cost}^{\left(n\right)}\left(x_{i},s_{i}\right)$,
% $y_{j}:=\mbox{arg max}_{x\in \mathbb{A}}\left(\bmu{n}(x, 1)+\tilde{\sigma}_n\left(x, z_{1:q}\right)W_{q}^{j}\right)$, and $\left\{ W_{q}^{k}\right\} _{k\geq1}$ are i.i.d standard normal random $q-$dimensional vectors. 

% Thus, first we observe that from Theorem 1, the expression of the stochastic gradient of $taKG+$ is:
% \begin{eqnarray*}
% G\left(Z\right) & =\\
% \frac{1}{r}\sum_{k=1}^{r}\frac{C\left(Z\right)\nabla\tilde{\sigma}_{n}\left(y_{k},Z\right)W_{q}^{k}}{C\left(Z\right)^{2}}\\
% -\frac{1}{r}\sum_{k=1}^{r}\frac{\nabla C\left(Z\right)}{C\left(Z\right)^{2}}\left(\mu^{\left(n\right)}\left(y_{k},1_{m}\right)+\tilde{\sigma}_{n}\left(y_{k},Z\right)W_{q}^{k}\right)
% \end{eqnarray*}
% where $C\left(Z\right):=\mbox{max}_{1\leq i\leq q}\mbox{cost}^{\left(n\right)}\left(x_{i},s_{i}\right)$,
% $y_{j}:=\mbox{arg max}_{x\in \mathbb{A}}\left(\bmu{n}(x, 1)+\tilde{\sigma}_n\left(x, z_{1:q}\right)W_{q}^{j}\right)$, and $\left\{ W_{q}^{k}\right\} _{k\geq1}$ are i.i.d standard normal random $q-$dimensional vectors. 

The theorem from \citet{kushner2003stochastic}, requires the following hypotheses:
\begin{enumerate}
\item $\epsilon_{t}\rightarrow0$, $\sum_{t=1}^{\infty}\epsilon_{t}=\infty$,
and $\sum_{t}\epsilon_{t}^{2}<\infty$.
\item $\sup_{t}E\left[\left|G\left(Z_{t}\right)\right|^{2}\right]<\infty$
\item There exist uniformly continuous functions $\left\{ \lambda_{t}\right\} _{t\geq0}$
of $Z$, and random vectors $\left\{ \beta_{t}\right\} _{t\geq0}$
, such that $\beta_{t}\rightarrow0$ almost surely and
\[
E_{n}\left[G\left(Z_{t}\right)\right]=\lambda_{t}\left(Z_{t}\right)+\beta_{t}.
\]
Furthermore, there exists a continuous function $\bar{\lambda}$,
such that for each $Z\in A^{q}$,
\[
\lim_{n}\left|\sum_{i=1}^{m\left(r_{m}+s\right)}\epsilon_{i}\left[\lambda_{i}\left(Z\right)-\bar{\lambda}\left(Z\right)\right]\right|=0
\]
for each $s\geq0$, where $m\left(r\right)$ is the unique value of
$k$ such that $t_{k}\leq t<t_{k+1}$, where $t_{0}=0$,$t_{k}=\sum_{i=0}^{k-1}\epsilon_{i}$.
\item There exists a continuously differentiable real-valued function $\phi$,
such that $\bar{\lambda}=-\nabla\phi$ and it is constant on each
connected subset of stationary points.
\item The constraint functions defining $\mathbb{A}$ are continuously differentiable.
\end{enumerate}
We now prove that our problem satisfy these hypotheses. (1) is true
by hypothesis of the lemma. 

Let's prove (2). We first assume that $r=1$, 
\begin{eqnarray*}
E\left[\left|\frac{1}{r}\sum_{k=1}^{r}\nabla\tilde{\sigma}_{n}\left(y_{k},Z\right)W_{q}^{k}\right|^{2}\right] & =\\
E\left[\left|\nabla\tilde{\sigma}_{n}\left(y_{1},Z\right)W_{q}^{1}\right|^{2}\right] & \leq\\
E\left[\left\Vert \nabla\tilde{\sigma}_{n}\left(y_{1},Z\right)\right\Vert ^{2}\left\Vert W_{q}^{1}\right\Vert ^{2}\right] & \leq\\
qL
\end{eqnarray*}
where $L:=\mbox{sup}_{x,z}\left\Vert \nabla\tilde{\sigma}_{n}\left(x,z\right)\right\Vert ^{2}$,
which is finite because the domain of the problem is compact and $\nabla\tilde{\sigma}_{n}\left(x,z\right)$
is continuous by Lemma 1. Since $C:=\mbox{cost}^{\left(n\right)}$ is continuously
differentiable bounded below by a constant $K$, thus we conclude
that the supremum over $Z$ of $E\left[\left|G\left(Z\right)\right|^{2}\right]$
is bounded. If $r>1$, $G\left(Z_{n}\right)$ is the average of i.i.d.
random vectors, whose squared norm expectation is finite, and so $\sup_{t}E\left[\left|G\left(Z_{t}\right)\right|^{2}\right]$
must be finite too.

We now prove (3). For each $t$, define 
\begin{eqnarray*}
\lambda_{t}\left(Z\right) & =\\
E\left[\frac{C\left(Z\right)\nabla\tilde{\sigma}_{n}\left(Y,Z\right)W_{q}}{C\left(Z\right)^{2}}\right]\\
-E\left[\frac{\nabla C\left(Z\right)}{C\left(Z\right)^{2}}\left(\mu_{n}\left(Y,1_{m}\right)+\tilde{\sigma}_{n}\left(Y,Z\right)W_{q}\right)\right]
\end{eqnarray*}
 where $Y=\mbox{arg max}_{x}\left(\mu_{n}\left(x\right)+\tilde{\sigma}_{n}\left(x,Z\right)W_{q}\right)$,
and $W_{q}$ is a standard normal random vector. Let's prove that
$\lambda_{t}$ is continuous. In the proof of Theorem 1, we show that
$\nabla\tilde{\sigma}_{n}\left(Y,Z\right)W_{q}$ is continuous in
$Z$. Furthermore,
\begin{eqnarray*}
\left\Vert \nabla\tilde{\sigma}_{n}\left(y_{1},Z\right)W_{q}^{1}\right\Vert  & \leq & \left\Vert \nabla\tilde{\sigma}_{n}\left(y_{1},Z\right)\right\Vert \left\Vert W_{q}\right\Vert \\
 & \leq & L\left\Vert W_{q}\right\Vert .
\end{eqnarray*}
Consequently $E\left[\nabla\tilde{\sigma}_{n}\left(Y,Z\right)W_{q}\right]$
is continuous by Corollary 5.7 of \citet{bartle}. In Theorem 1, we also show
that $E\left[\left(\mu_{n}\left(Y,1_{m}\right)+\tilde{\sigma}_{n}\left(Y,Z\right)W_{q}\right)\right]$
is continuous in $Z$. Since $C$ is continuously differentiable,
we conclude that $\lambda_{t}$ is continuous. By defining $\beta_{t}=0$
for all $t$, and $\bar{\lambda}=\lambda_{1}$, we conclude the proof
of (3).

Finally, define $\phi\left(Z\right)=-E\left[\frac{\mu_{n}\left(Y,1_{m}\right)+\tilde{\sigma}_{n}\left(Y,Z\right)W_{q}}{C\left(Z\right)}\right]$.
Observe that in Lemma 2, we show that we can interchange the expectation
and the gradient in $E\left[\nabla\left(\mu_{n}\left(Y\right)+\tilde{\sigma}_{n}\left(Y,Z\right)W_{q}\right)\right]$,
and so $\lambda_{m}\left(Z\right)=-\nabla\phi\left(Z\right).$ In
a connected subset of stationary points, we have that $\lambda_{m}\left(Z\right)=0$,
and so $\phi\left(Z\right)$ is constant. This ends the proof of the
theorem.
\endproof

\begin{proof}[proof of Proposition 1]
Since
\begin{eqnarray*}
VOI_{n}(x, s) &:=& \\
\E_n[\mu^*(x, 1) - \min_{x'} (\mu_n(x') + C_{n}(x', (x, s)) W]
\end{eqnarray*}
where $W$ is a standard normal random variable. By Jensen's inequality, we have
\begin{eqnarray*}
VOI_{n}(x, s) &:=& \E_n[\mu^*(x, 1) - \min_{x'} (u_{n}\left(x',s,W\right))]\\
&\ge& \mu^*(x, 1) - \min_{x'} \E_n(u_{n}\left(x',s,W\right)) = 0.
\end{eqnarray*}
where $u_{n}\left(x,s,W\right):=\mu^n(x', 1) + C_{n}((x', 1), (x, s)) W)$.
%\begin{eqnarray*}
%VOI(x, s) &:=& \E_n[\mu^*(x, 1) - \min_{x'} (\mu^n(x', 1) + C^{n}((x', 1), (x, s)) W)]\\
%&\ge& \mu^*(x, 1) - \min_{x'} \E_n(\mu^n(x') + C^{n}((x', 1), (x, s)) W) = 0.
%\end{eqnarray*}
The inequality becomes equal only if $\min_{x'} (\mu_n(x') + C_{n}(x', (x, s)) W$ is a linear function of $W$ for any fixed $(x, s)$, i.e the argmin for the inner optimization function doesn't change as we vary $W$, which is not true if $K_{n}((x', 1), (x, s))>0$ i.e. evaluating at $(x, s)$ provides value to determine the argmin of the surface $(x, 1)$.
\end{proof}

\begin{proof}[proof of Proposition 2]
The proof follows a very similar argument than the previous proof. By Jensen's inequality, we have that

\begin{eqnarray*}
\E_{n}\left[\min_{x'}\E_{n}\left[g(x',1)\mid\Y(x,\S)\right]\right] & \geq\\
\E_{n}\left[\min_{x'}\E_{n}\left[g(x',1)\mid\Y(x,\S\bigcup C(\S))\right]\right]
\end{eqnarray*}

The inequality becomes equal only if the argmin for the inner optimization function doesn't change as we vary the normal random vector, which is not true under our assumptions.


\end{proof}

\subsection{GPs for Hyperparameter Optimization}
\label{sect:gp-ho}
In the context of hyperparameter optimization with two continuous fidelities, i.e. the number of training iterations ($s_{(1)}$) and the amount of training data ($s_{(2)}$), we set the kernel function of the GP as
\begin{eqnarray*}
K_{0}(z, \tilde z) = K(x, \tilde x) \times K_1(s_{(1)}, \tilde s_{(1)}) \times K_2(s_{(2)},, \tilde s_{(2)}),
\end{eqnarray*}
where $K(\cdot, \cdot)$ is a square-exponential kernel. If we assume that the learning curve looks like
\begin{eqnarray}
g(x, s) = h(x) \times \left(\beta_0 + \beta_1\exp{(-\lambda s_{(1)})}\right) \times l(s_{(2)}),
\label{eqn:shape}
\end{eqnarray}
then inspired by \cite{swersky2014freeze}, we set the kernel $K_1(\cdot, \cdot)$ as
\begin{eqnarray*}
K_1(s_{(1)}, \tilde s_{(1)}) = \left( w + \frac{\beta^{\alpha}}{(s_{(1)} + \tilde s_{(1)} + \beta^{\alpha})}\right),
\end{eqnarray*}
where $w, \beta, \alpha > 0$ are hyperparameters. We add an intercept $w$ compared to the kernel in \cite{swersky2014freeze} to model the fact that the loss will not diminish. We assume that the kernel $K_2(\cdot, \cdot)$ has the form
\begin{eqnarray*}
K_2(s_{(2)}, \tilde s_{(2)}) = \left(c + (1-s_{(2)})^{(1+\delta)}(1-\tilde s_{(2)})^{(1+\delta)}\right),
\end{eqnarray*}
where $c, \delta > 0$ are hyperparameters.

All the hyperparameters can be treated in a Bayesian way as proposed in \citet{snoek2012practical}.

\subsection{Additional experimental details}
\label{sect:addition}
\subsubsection{Synthetic experiments}
Here we define in detail the synthetic test functions on which we perform numerical experiments
The test functions are: 
\begin{eqnarray*}
&&\text{augmented-Branin}(x, s) \\
&=& \left(x_2 - \left(\frac{5.1}{4 \pi^2} - 0.1*(1-s_1)\right) x_1^2 + \frac{5}{\pi} x_1 - 6\right)^2 \\
&& \quad + 10*\left(1-\frac{1}{8\pi}\right)\cos(x_1) + 10\\
&&\text{augmented-Hartmann}(x, s) \\
&=& \left(\alpha_1-0.1*(1-s_1)\right) \exp{\left(-\sum_{j=1}^d A_{ij} (x_j - P_{1j})^2\right)}\\
&&\quad + \sum_{i=2}^{4} \alpha_i \exp{\left(-\sum_{j=1}^d A_{ij} (x_j - P_{ij})^2\right)}\\
&&\text{augmented-Rosenbrock}(x, s) \\
& = & \sum_{i=1}^2 \left(100*(x_{i+1}-x_i^2 + 0.1*(1-s_1))^2\right.\\
& & + \left.\left(x_i-1+0.1*(1-s_2)^2\right)^2\right).
\end{eqnarray*} 

\subsubsection{Real-world experiments}
The range of search domain for feedforward NN experiments: the learning rate in $[10^{-6}, 10^0]$, dropout rate in $[0, 1]$, batch size in $[2^5, 2^{10}]$ and the number of units at each layer in $[100, 1000]$.

The range of search domain for CNN experiments: the learning rate in $[10^{-6}, 1.0]$, batch size $[2^5, 2^{10}]$, and number of filters in each convolutional block in $[2^5, 2^9]$. 



\bibliography{supplement}



\end{document}
